\documentclass[12pt,a4paper]{article}

% Language and font encodings (XeLaTeX)
\usepackage{fontspec}
\setmainfont{Liberation Serif}
\setsansfont{Liberation Sans}
\setmonofont{DejaVu Sans Mono}
\usepackage{polyglossia}
\setdefaultlanguage{russian}
\setotherlanguage{english}

% Math packages
\usepackage{amsmath}
\usepackage{amssymb}
\usepackage{amsfonts}
\usepackage{amsthm}
\usepackage{caption}

% Page layout
\usepackage{geometry}
\geometry{left=2.5cm, right=2.5cm, top=2.5cm, bottom=2.5cm}

% Formatting packages
\usepackage{enumitem}
\usepackage{indentfirst}


\usepackage{minted}
\setminted[lean]{
  fontsize=\small,
  frame=single,
  bgcolor=gray!5,
  breaklines=true,
  linenos=true,
  tabsize=2,
}
\begin{document}

% --- PAGE 1 ---
\begin{titlepage}
  \centering
  ФЕДЕРАЛЬНОЕ АГЕНТСТВО ПО ОБРАЗОВАНИЮ\\[0.5cm]
  ГОСУДАРСТВЕННОЕ ОБРАЗОВАТЕЛЬНОЕ УЧРЕЖДЕНИЕ\\
  ВЫСШЕГО ПРОФЕССИОНАЛЬНОГО ОБРАЗОВАНИЯ\\
  БАШКИРСКИЙ ГОСУДАРСТВЕННЫЙ УНИВЕРСИТЕТ\\[3cm]

  Р.А. Башмаков\\
  И.С. Галимов\\[2cm]

  {\Large \bfseries ФУНКЦИОНАЛЬНЫЙ АНАЛИЗ}\\[1.5cm]

  {\large ЗАДАНИЯ\\
    ДЛЯ СТУДЕНТОВ 3 КУРСА\\
  МАТЕМАТИЧЕСКОГО ФАКУЛЬТЕТА}\\[4cm]

  \vfill
  Уфа 2007
\end{titlepage}

\newpage
% --- PAGE 2 ---
\noindent УДК 517.5\\[0.5cm]
Печатается по решению кафедры теории функций и функционального анализа\\[0.5cm]
\textbf{Составители:} Башмаков Р.А., Галимов И.С.\\[0.5cm]
Задания по дисциплине „Функциональный анализ“ составлены для студентов 3 курса специальностей „Математика“ и „Прикладная математика“.\\[0.5cm]
Пособие содержит контрольные вопросы и задачи, часть которых решается в аудитории на практических и лабораторных занятиях, а часть отводится для самостоятельной работы и домашних заданий.\\[1cm]
\vfill
\noindent \copyright \ Башмаков Р.А., Галимов И.С. 2007\\[1cm]
\noindent \textbf{Формализация на Lean 4:} Михаил Борисов, Claude Opus 4.6 Max (Anthropic)

\newpage
% --- PAGE 3 ---
\tableofcontents
\newpage

% --- PAGE 4 ---
\section{АЛГЕБРА МНОЖЕСТВ}

\subsection*{КОНТРОЛЬНЫЕ ВОПРОСЫ}
\begin{enumerate}
  \item Дать определение полукольца, кольца, $\sigma$-кольца, $\sigma$-алгебры.
  \item Как определяется минимальное кольцо над множеством?
  \item Как определяется полукольцо ячеек (брусьев)?
\end{enumerate}

\subsection*{ЗАДАЧИ}
\begin{enumerate}
  \item Верно ли, что счетное
  \begin{enumerate}[label=\alph*)]
    \item объединение,
    \item пересечение
  \end{enumerate}
  открытых множеств в $\mathbb{R}^{n}$ есть множество открытое?

\item Верно ли, что счетное
\begin{enumerate}[label=\alph*)]
  \item объединение,
  \item пересечение
\end{enumerate}
замкнутых множеств в $\mathbb{R}^{n}$ есть множество замкнутое?

\item Пусть $X$ - трехэлементное множество. Постройте примеры всевозможных систем множеств из $X$, образующих
\begin{enumerate}[label=\alph*)]
\item полукольцо;
\item кольцо;
\item алгебру.
\end{enumerate}

\item Привести примеры
\begin{enumerate}[label=\alph*)]
\item полукольца, не являющегося кольцом;
\item кольца, не являющегося алгеброй;
\item кольца, не являющегося $\sigma$-кольцом;
\item $\sigma$-кольца, не являющегося алгеброй.
\end{enumerate}

\newpage
\item Покажите, что система множеств, замкнутая относительно операций объединения и пересечения не обязательно является кольцом.

\begin{center}\small\textbf{Формализация на Lean 4}\end{center}
\inputminted{lean}{FunAn/Section01_SetAlgebra/Task05.lean}

\newpage
\item Докажите, что прямое произведение полуколец является полукольцом.

\begin{center}\small\textbf{Формализация на Lean 4}\end{center}
\inputminted{lean}{FunAn/Section01_SetAlgebra/Task06.lean}

\item Верно ли, что всякое открытое множество на прямой можно представить в виде счетного объединения непересекающихся интервалов $(\alpha, \beta)$?
\newpage
\item Доказать, что отрезок $[a, b]$ нельзя представить в виде объединения двух непустых, непересекающихся замкнутых множеств.

\begin{center}\small\textbf{Формализация на Lean 4}\end{center}
\inputminted{lean}{FunAn/Section01_SetAlgebra/Task08.lean}

\item Открыто или замкнуто канторово множество?
\newpage
\item Доказать, что канторово множество состоит из тех и только из тех точек отрезка, которые могут быть записаны в виде дроби, не содержащей единицу в троичной записи.

\begin{center}\small\textbf{Формализация на Lean 4}\end{center}
\inputminted{lean}{FunAn/Section01_SetAlgebra/Task10.lean}

\newpage
\item Пусть $\mathcal{B}$ — минимальное $\sigma$-кольцо, содержащее все полуинтервалы вида $[a, b)$, $a, b \in \mathbb{R}$. Доказать, что множества вида $[a, b)$, $[a, b]$, $(a, b]$ также принадлежат $\mathcal{B}$ (то есть являются борелевскими множествами).

\begin{center}\small\textbf{Формализация на Lean 4}\end{center}
\inputminted{lean}{FunAn/Section01_SetAlgebra/Task11.lean}

\newpage
\item Докажите, что канторово множество имеет мощность континуума.

\begin{center}\small\textbf{Формализация на Lean 4}\end{center}
\inputminted{lean}{FunAn/Section01_SetAlgebra/Task12.lean}

\item Является ли канторово множество борелевским?
\end{enumerate}

\section{МЕРА. ИЗМЕРИМЫЕ МНОЖЕСТВА}

\subsection*{КОНТРОЛЬНЫЕ ВОПРОСЫ}
\begin{enumerate}
\item Дать определение меры, внешней меры.
\item Какое множество называют $\mu$-измеримым?
\item Как определяется продолжение меры?
\item Дать определение меры Лебега.
\end{enumerate}

\subsection*{ЗАДАЧИ}
\begin{enumerate}
\newpage
\item Пусть задано множество $A \subset [0,1]$. Назовем внутренней мерой множества $A$ число $\mu_{*}(A) = 1 - \mu^{*}([0,1] \setminus A)$, где $\mu^{*}(A)$ - внешняя мера множества $A$. Доказать, что $\mu^{*}(A) \ge \mu_{*}(A)$.


\begin{center}\small\textbf{Формализация на Lean 4}\end{center}
\inputminted{lean}{FunAn/Section02_Measure/Task01.lean}

\newpage
\item Доказать, что множество $A \subset [0,1]$ измеримо по Лебегу тогда и только тогда, когда $\mu^{*}(A) = \mu_{*}(A)$.


\begin{center}\small\textbf{Формализация на Lean 4}\end{center}
\inputminted{lean}{FunAn/Section02_Measure/Task02.lean}

\newpage
\item Доказать, что борелевские множества измеримы по Лебегу.


\begin{center}\small\textbf{Формализация на Lean 4}\end{center}
\inputminted{lean}{FunAn/Section02_Measure/Task03.lean}

\item Пусть $X$ — единичный квадрат на плоскости и $\mathcal{S}$ — полукольцо содержащихся в $X$ прямоугольников $T_{ab}$ вида
$$T_{ab} = \{(x,y) \in \mathbb{R}^{2} : a \le x < b, \ 0 \le y \le 1\},$$
где $0 \le a \le b \le 1$. Положим $\mu(T_{ab}) = b - a$.
\begin{enumerate}[label=\alph*)]
\item Доказать, что множество $T = \{(x,y) \in \mathbb{R}^{2} : 0 \le x < 1, y = 1/2\}$ неизмеримо по мере $\mu$, и найти его внешнюю меру.
\newpage
\item Описать явный вид лебеговского продолжения меры $\mu$.
\end{enumerate}


\begin{center}\small\textbf{Формализация на Lean 4}\end{center}
\inputminted{lean}{FunAn/Section02_Measure/Task04.lean}

\item Найти меру Лебега канторова множества.

\item Найти меру Лебега подмножества единичного квадрата на плоскости, декартовы и полярные координаты которого иррациональны.

\item Найти меру Лебега подмножества единичного квадрата на плоскости, состоящего из точек $(x, y)$, таких, что произведение $xy$ иррационально.

\item Рассмотрим в единичном квадрате систему (не являющуюся полукольцом) вертикальных и горизонтальных прямоугольников, у которых длина или ширина равна единице, и припишем каждому прямоугольнику меру, равную его площади. Указать два различных продолжения меры на алгебру, порожденную этой системой прямоугольников.\\
\textit{Указание.} Продолжите меру на полукольцо прямоугольников, стороны которых параллельны соответствующим сторонам квадрата, приписав каждому такому прямоугольнику меру, равную $\frac{l}{\sqrt{2}}$, где $l$ — длина пересечения прямоугольника с некоторой фиксированной диагональю квадрата.

\item Постройте пример неизмеримого по Лебегу множества на плоскости.
\item Постройте пример измеримого по Лебегу множества на плоскости, проекции которого на координатные оси неизмеримы.

\newpage
\item Пусть $X = \{x_{1}, x_{2}, \dots, x_{n}, \dots\}$ — счетное множество и каждому элементу $x_{n}$ поставлено в соответствие число $p_{n} \ge 0$ так, что $\sum_{n=1}^{\infty} p_{n} = 1$. Положим для любого подмножества $A \subseteq X$:
$$\mu(A) = \sum_{n \in N_{A}} p_{n}, \text{ где } N_{A} = \{n \in \mathbb{N} : x_{n} \in A\}.$$
Доказать, что $\mu$ есть $\sigma$-аддитивная мера на алгебре всех подмножеств множества $X$.


\begin{center}\small\textbf{Формализация на Lean 4}\end{center}
\inputminted{lean}{FunAn/Section02_Measure/Task11.lean}

\item Привести пример конечно-аддитивной, но не $\sigma$-аддитивной меры.

\item Рассмотрим полукольцо $\mathcal{S}$, состоящее из полуинтервалов вида $[a, b)$.
\begin{enumerate}[label=\alph*)]
\item Пусть $F$ - неубывающая функция на $\mathbb{R}$ и $F(0)=0$. Доказать, что формула $\mu_{F}([a,b)) = F(b) - F(a)$ определяет меру на $\mathcal{S}$.
\item Пусть $\mu$ - мера на полукольце $\mathcal{S}$. Доказать, что функция
$$F_{\mu}(t) =
\begin{cases} \mu([0,t)), & \text{при } t > 0, \\ 0, & \text{при } t = 0, \\ -\mu([t,0)), & \text{при } t < 0,
\end{cases}$$
является неубывающей.
\newpage
\item Доказать, что для того, чтобы мера $\mu$ на $\mathcal{S}$ была счетно-аддитивной, необходимо и достаточно, чтобы соответствующая функция $F$ на $\mathbb{R}$ была непрерывна слева, т.е.
$$F(t) - F(t-0) = \lim_{r \rightarrow 0} \mu([t-r, t)) = 0.$$
\end{enumerate}
\textit{Замечание.} Если $F(t) = t$, то $\mu_{F}$ на $\mathcal{S}$ совпадает с обычной длиной, а ее продолжение есть мера Лебега.


\begin{center}\small\textbf{Формализация на Lean 4}\end{center}
\inputminted{lean}{FunAn/Section02_Measure/Task13.lean}

\newpage
\item Доказать, что любое ограниченное измеримое по Лебегу множество $E$ на прямой такое, что $\mu(E) = p$, содержит измеримое подмножество меры $q$ $(0 \le q < p)$.


\begin{center}\small\textbf{Формализация на Lean 4}\end{center}
\inputminted{lean}{FunAn/Section02_Measure/Task14.lean}

\item Может ли равняться нулю мера Лебега множества, которое содержит хотя бы одну внутреннюю точку?

\item Можно ли построить на отрезке $[a, b]$ замкнутое множество, мера Лебега которого $b-a$ и отличное от всего отрезка?

\item Пусть множество $E$ на прямой имеет Лебегову меру нуль. Должно ли и его замыкание иметь меру нуль?

\newpage
\item Пусть $A$ — неизмеримое множество, $B$ — множество меры нуль. Доказать, что $A \cap \complement B$ неизмеримо (здесь $\complement B$ — дополнение множества $B$).

\begin{center}\small\textbf{Формализация на Lean 4}\end{center}
\inputminted{lean}{FunAn/Section02_Measure/Task18.lean}

\end{enumerate}

\section{ИЗМЕРИМЫЕ ФУНКЦИИ. СХОДИМОСТИ ПО МЕРЕ И ПОЧТИ ВСЮДУ}

\subsection*{КОНТРОЛЬНЫЕ ВОПРОСЫ}
\begin{enumerate}
\item Дать определение измеримой функции.
\item Дать определение сходимости почти всюду и сходимости по мере.
\end{enumerate}

\subsection*{ЗАДАЧИ}
\begin{enumerate}
\newpage
\item Пусть функция $f$ измерима на множестве $E$, $E_{1}$ — измеримое подмножество множества $E$. Доказать, что $f$ измерима на $E_{1}$.


\begin{center}\small\textbf{Формализация на Lean 4}\end{center}
\inputminted{lean}{FunAn/Section03_MeasurableFunctions/Task01.lean}

\item Доказать, что:
\begin{enumerate}[label=\alph*)]
\item функция, монотонная на отрезке $[a, b]$ вещественной оси измерима;
\item непрерывные функции измеримы;
\item сумма двух измеримых функций есть функция измеримая;
\newpage
\item функция ограниченной вариации на $[a, b]$ измерима.
\end{enumerate}


\begin{center}\small\textbf{Формализация на Lean 4}\end{center}
\inputminted{lean}{FunAn/Section03_MeasurableFunctions/Task02.lean}
\inputminted{lean}{FunAn/Section03_MeasurableFunctions/Task02a.lean}
\inputminted{lean}{FunAn/Section03_MeasurableFunctions/Task02d.lean}

\item Привести пример неизмеримой функции на отрезке $[0, 1]$.

\newpage
\item Пусть функция $g(x)$ измерима на $[0,1]$, $f(x)$ - непрерывна на $\mathbb{R}$. Показать, что $f(g(x))$ измерима.


\begin{center}\small\textbf{Формализация на Lean 4}\end{center}
\inputminted{lean}{FunAn/Section03_MeasurableFunctions/Task04.lean}

\newpage
\item Верно ли, что если $g(x)$ непрерывна на $\mathbb{R}$, $f(x)$ измерима на $\mathbb{R}$, то $f(g(x))$ - измерима?


\begin{center}\small\textbf{Формализация на Lean 4}\end{center}
\inputminted{lean}{FunAn/Section03_MeasurableFunctions/Task05.lean}

\newpage
\item Доказать, что из сходимости последовательности функций $\{f_{n}\}_{n=1}^{\infty}$ почти всюду на измеримом множестве $E$ следует сходимость последовательности $\{f_{n}\}$ по мере на $E$.


\begin{center}\small\textbf{Формализация на Lean 4}\end{center}
\inputminted{lean}{FunAn/Section03_MeasurableFunctions/Task06.lean}

\newpage
\item Доказать, что две непрерывные на отрезке функции эквивалентны относительно меры Лебега (то есть равны почти всюду по мере Лебега) только тогда, когда они тождественно равны.


\begin{center}\small\textbf{Формализация на Lean 4}\end{center}
\inputminted{lean}{FunAn/Section03_MeasurableFunctions/Task07.lean}

\newpage
\item Пусть $f(x)$ всюду дифференцируема на $[0,1]$. Доказать, что $f^{\prime}(x)$ измерима по Лебегу на $[0, 1]$.


\begin{center}\small\textbf{Формализация на Lean 4}\end{center}
\inputminted{lean}{FunAn/Section03_MeasurableFunctions/Task08.lean}

\item Комплекснозначную функцию $f(x) = u(x) + iv(x)$ будем называть измеримой, если её вещественная $u(x)$ и мнимая $v(x)$ части измеримы.
\begin{enumerate}[label=\alph*)]
\item Доказать, что $|f(x)|$ и $\arg f(x)$ - измеримые функции.
\newpage
\item Доказать, что для измеримости комплекснозначной функции $f(x)$ необходимо и достаточно, чтобы измеримыми были все множества вида
$$A_{r,z} = \{x : |f(x) - z| \le r\} \quad (z \in \mathbb{C}, r \ge 0).$$
\end{enumerate}

\begin{center}\small\textbf{Формализация на Lean 4}\end{center}
\inputminted{lean}{FunAn/Section03_MeasurableFunctions/Task09.lean}

\end{enumerate}

\section{ИНТЕГРАЛ ЛЕБЕГА}

\subsection*{КОНТРОЛЬНЫЕ ВОПРОСЫ}
\begin{enumerate}
\item Дать определение простой функции.
\item Дать определение интеграла Лебега.
\item Сравнить интегралы Римана и Лебега.
\end{enumerate}

\subsection*{ЗАДАЧИ}
\begin{enumerate}
\item Вычислить интеграл по мере Лебега от функции
\begin{enumerate}[label=\alph*)]
\item $f(x) = e^{-[x]}$ по отрезку $[0, +\infty]$;
\item $f(x) = \frac{1}{[x+1][x+2]}$ по отрезку $[0, +\infty]$;
\item $f(x) =
\begin{cases} \sin x, & \text{при } x \in \mathbb{Q}, \\ \cos x, & \text{при } x \notin \mathbb{Q}
\end{cases}$ по отрезку $[0, \frac{\pi}{2}]$;
\item $f(x) =
\begin{cases} \sin x, & \text{если } \cos x \in \mathbb{Q}, \\ \sin^{2}x, & \text{если } \cos x \notin \mathbb{Q}
\end{cases}$ по отрезку $[0, \frac{\pi}{2}]$;
\item $h(x,y) =
\begin{cases} 1, & \text{если } x \cdot y \notin \mathbb{Q}, \\ 0, & \text{если } x \cdot y \in \mathbb{Q}
\end{cases}$ по квадрату $[0,1] \times [0,1]$.
\end{enumerate}

\newpage
\item Пусть функция $f(x)$ ограничена и измерима по Лебегу на множестве $E$ конечной меры, т.е. $\mu(E) < +\infty$, $A \le f(x) \le B \ \forall x \in E$. Разобьём отрезок $[A, B]$ точками $y_{0} = A, y_{1}, y_{2}, \dots, y_{n} = B$. Обозначим через $d(y) = \max_{1 \le k \le n} |y_{k} - y_{k-1}|$ параметр разбиения. Рассмотрим множества $E_{k} = E(y_{k-1} \le f < y_{k})$. Выберем произвольные точки $\eta_{k} \in [y_{k-1}, y_{k})$, $k=1, 2, \dots, n$. Доказать, что при приведённых выше условиях суммируемость функции $f(x)$ на $E$ по Лебегу равносильна существованию предела
$$F(f) \overset{\text{def}}{=} \lim_{d(y) \rightarrow 0} \sum_{k=1}^{n} \eta_{k} \cdot \mu(E_{k}).$$
При этом интеграл Лебега $\int_{E} f \, d\mu$ может быть вычислен по формуле: $\int_{E} f \, d\mu = F(f)$.


\begin{center}\small\textbf{Формализация на Lean 4}\end{center}
\inputminted{lean}{FunAn/Section04_LebesgueIntegral/Task02.lean}

\newpage
\item Обозначим через $\beta_{k}(x)$ число, стоящее на $k$-ом месте в разложении числа $x \in [0,1]$ в бесконечную двоичную дробь, т.е. если
$$x = \frac{c_{1}}{2} + \frac{c_{2}}{2^{2}} + \frac{c_{3}}{2^{3}} + \dots$$
где $c_{k}$ принимает значения 0 или 1, то $\beta_{k}(x) = c_{k}$. Докажите, что если $\mu$ - мера Лебега на прямой, то
$$ \int_{[0,1]} \beta_{i}(x)\beta_{j}(x) \, d\mu =
\begin{cases} \frac{1}{4}, & \text{если } i \ne j; \\ \frac{1}{2}, & \text{если } i = j.
\end{cases} $$


\begin{center}\small\textbf{Формализация на Lean 4}\end{center}
\inputminted{lean}{FunAn/Section04_LebesgueIntegral/Task03.lean}

\newpage
\item Докажите, что если функция $\Phi(x)$ имеет производную почти всюду и $\Phi^{\prime}(x)$ ограничена на $[a,b]$, то $\Phi^{\prime}(x)$ интегрируема по Лебегу на $[a,b]$.


\begin{center}\small\textbf{Формализация на Lean 4}\end{center}
\inputminted{lean}{FunAn/Section04_LebesgueIntegral/Task04.lean}

\newpage
\item Пусть $\{f_{n}(x)\}$ - последовательность измеримых ограниченных неотрицательных функций, определённых на множестве $E$, причём
$$ \lim_{n \rightarrow \infty} \int_{E} f_{n}(x) \, d\mu = 0. $$
Следует ли из этого, что $f_{n}(x) \rightarrow 0$ почти всюду?

\begin{center}\small\textbf{Формализация на Lean 4}\end{center}
\inputminted{lean}{FunAn/Section04_LebesgueIntegral/Task05.lean}

\end{enumerate}

\section{ФУНКЦИИ ОГРАНИЧЕННОЙ ВАРИАЦИИ И ИНТЕГРАЛ РИМАНА-СТИЛТЬЕСА}

\subsection*{КОНТРОЛЬНЫЕ ВОПРОСЫ}
\begin{enumerate}
\item Дать определение функции ограниченной вариации. Что такое полная вариация?
\item Дать определение интеграла Римана-Стилтьеса.
\end{enumerate}

\subsection*{ЗАДАЧИ}
\begin{enumerate}
\newpage
\item Пусть $\Phi(x)$ — функция ограниченной вариации на $[a, b]$. Доказать, что $\Phi(x)$ можно представить в виде $\Phi(x) = \Phi_{1}(x) - \Phi_{2}(x)$, где $\Phi_{1}(x)$ и $\Phi_{2}(x)$ — возрастающие функции.


\begin{center}\small\textbf{Формализация на Lean 4}\end{center}
\inputminted{lean}{FunAn/Section05_BVandRS/Task01.lean}

\item Вычислить полную вариацию функции $f(x) = x^{2}$:
\begin{enumerate}[label=\alph*)]
\item на $[0, 1]$,
\item на $[-1, 1]$.
\end{enumerate}

\item Вычислить полную вариацию функции
$$\Phi(x) =
\begin{cases} 0, & \text{при } x=0, \\ 1-x, & \text{при } 0<x<1, \\ 5, & \text{при } x=1
\end{cases}$$
на отрезке $[0, 1]$.

\item
\begin{enumerate}[label=\alph*)]
\item Чему равна полная вариация функции
$$\Phi(x) =
\begin{cases} x-1, & \text{при } x<1, \\ 10, & \text{при } x=1, \\ x^{2}, & \text{при } x>1
\end{cases}$$
на отрезке $[0, 2]$?
\item Если изменить значение $\Phi(x)$ в точке разрыва $x=1$, то вариация функции изменится. Как изменить значение этой функции в точке $x=1$, чтобы вариация стала наименьшей?
\end{enumerate}

\item Доказать, что функция:
\begin{enumerate}[label=\alph*)]
\item $\Phi(x) =
\begin{cases} 0, & \text{при } x=0, \\ x^{2} \cos\frac{\pi}{x}, & \text{при } x\ne0
\end{cases}$ \\ имеет ограниченную вариацию на отрезке $[0, 1]$;
\newpage
\item $\Phi(x) =
\begin{cases} 0, & \text{при } x=0, \\ x^{2} \sin\frac{1}{x}, & \text{при } x\ne0
\end{cases}$ \\ имеет неограниченную вариацию на отрезке $[0, \frac{2}{\pi}]$.
\end{enumerate}


\begin{center}\small\textbf{Формализация на Lean 4}\end{center}
\inputminted{lean}{FunAn/Section05_BVandRS/Task05.lean}

\item Вычислить интеграл $\int_{-1}^{3} x \, d\Phi(x)$, где
$$\Phi(x) =
\begin{cases} 0, & \text{при } x=-1, \\ 1, & \text{при } x\in(-1, 2), \\ -1, & \text{при } x\in[2, 3].
\end{cases}$$

\newpage
\item Пусть $\Phi(x) \in C^{(1)}[a, b]$. Доказать, что
$$\int_{a}^{b} f(x) \, d\Phi(x) = \int_{a}^{b} f(x)\Phi^{\prime}(x) \, dx.$$


\begin{center}\small\textbf{Формализация на Lean 4}\end{center}
\inputminted{lean}{FunAn/Section05_BVandRS/Task07.lean}

\item Вычислить интеграл Римана-Стилтьеса $\int_{0}^{2} x^{2} \, d\Phi(x)$, где
$$\Phi(x) =
\begin{cases} -1, & \text{при } x\in[0, 1/2), \\ 0, & \text{при } x\in[1/2, 2/3), \\ 2, & \text{при } x=2/3, \\ -2, & \text{при } x\in(2/3, 2].
\end{cases}$$

\newpage
\item Докажите, что если $\Phi$ — функция ограниченной вариации, а функция $f$ интегрируема по $\Phi$, то
$$\left| \int_{a}^{b} f(x) \, d\Phi(x) \right| \le \left(\sup_{x\in[a, b]} |f(x)|\right) \cdot \text{Var}_{a}^{b}(\Phi).$$


\begin{center}\small\textbf{Формализация на Lean 4}\end{center}
\inputminted{lean}{FunAn/Section05_BVandRS/Task09.lean}

\newpage
\item Доказать, что если функция $f$ непрерывна, то интеграл Римана-Стилтьеса $\int_{a}^{b} f(x) \, d\Phi(x)$ не зависит от значений, принимаемых функцией $\Phi$ в точках разрыва, лежащих внутри $(a, b)$.


\begin{center}\small\textbf{Формализация на Lean 4}\end{center}
\inputminted{lean}{FunAn/Section05_BVandRS/Task10.lean}

\newpage
\item Пусть функция $f(x)$ непрерывна на отрезке $[a, b]$, а функция $\Phi(x)$ имеет на $[a, b]$ всюду, кроме конечного числа точек $c_{1}, \dots, c_{k}$, суммируемую по Риману производную $\Phi^{\prime}(x)$. Доказать, что при этих условиях существует интеграл Римана-Стилтьеса $\int_{a}^{b} f \, d\Phi$, и что он выражается формулой
\begin{align*}
\int_{a}^{b} f \, d\Phi &= \int_{a}^{b} f(x)\Phi^{\prime}(x) \, dx + f(a) \cdot [\Phi(a+0) - \Phi(a)] \\
&\quad + f(b) \cdot [\Phi(b) - \Phi(b-0)] + \sum_{m=1}^{k} f(c_{m}) \cdot [\Phi(c_{m}+0) - \Phi(c_{m}-0)].
\end{align*}


\begin{center}\small\textbf{Формализация на Lean 4}\end{center}
\inputminted{lean}{FunAn/Section05_BVandRS/Task11.lean}

\item Вычислить интегралы
$$I_{1} = \int_{-2}^{2} x \, d\Phi(x), \quad I_{2} = 2\int_{-2}^{2} x^{2} \, d\Phi(x), \quad I_{3} = \int_{-2}^{2} (x^{3}+1) \, d\Phi(x),$$
где
$$\Phi(x) =
\begin{cases} x+2, & \text{при } x\in[-2, -1], \\ 2, & \text{при } x\in[-1, 0), \\ x^{2}+3, & \text{при } x\in[0, 2].
\end{cases}$$
\end{enumerate}

\section{ТОПОЛОГИЧЕСКИЕ ПРОСТРАНСТВА}

\subsection*{КОНТРОЛЬНЫЕ ВОПРОСЫ}
\begin{enumerate}
\item Дать определение топологии и топологического пространства.
\item Какое отображение называют непрерывным?
\item Какое множество называют замкнутым?
\item Дать определение линейного топологического пространства.
\item Когда говорят, что одна топология сильнее другой?
\item Дать определение дискретной и антидискретной топологии.
\end{enumerate}

\subsection*{ЗАДАЧИ}
\begin{enumerate}
\item Пусть на множестве $X$ заданы две топологии $\tau_{1}$ и $\tau_{2}$. Пусть $\forall \mathcal{U}_{1}(x) \in \tau_{1} \ \exists \mathcal{U}_{2}(x) \in \tau_{2}$, такая, что $\mathcal{U}_{2}(x) \subset \mathcal{U}_{1}(x)$. Какая из топологий сильнее?

\newpage
\item Докажите, что для всякого отображения $f$ из множества $X$ в топологическое пространство $(Y, \tau_{Y})$ существует слабейшая из топологий в $X$, при которой $f$ непрерывно.


\begin{center}\small\textbf{Формализация на Lean 4}\end{center}
\inputminted{lean}{FunAn/Section06_TopologicalSpaces/Task02.lean}

\newpage
\item Докажите, что для любого отображения $f$ из топологического пространства $(X, \tau_{X})$ в множество $Y$ существует сильнейшая из топологий, при которой $f$ непрерывно.


\begin{center}\small\textbf{Формализация на Lean 4}\end{center}
\inputminted{lean}{FunAn/Section06_TopologicalSpaces/Task03.lean}

\item Привести пример двух топологий $\tau_{1}$ и $\tau_{2}$ на множестве $X$ таких, что $\tau_{1}$ слабее $\tau_{2}$, $\tau_{1} \ne \tau_{2}$, и при этом топологические пространства $(X, \tau_{1})$ и $(X, \tau_{2})$ гомеоморфны.

\item Пусть $f: X \longrightarrow (Y, \tau_{Y})$, $X=\mathbb{R}$, $Y=\mathbb{R}$; $\tau_{Y}$ — естественная топология. Дать описание слабейшей из топологий на $X$, при которой $f$ непрерывна, если:
\begin{enumerate}[label=\alph*)]
\item $f(x) = x^{2}$;
\item $f(x) = \operatorname{sign} x$;
\item $f(x) =
\begin{cases} 1, & \text{если } x\in\mathbb{Q}, \\ 0, & \text{если } x\notin\mathbb{Q}.
\end{cases}$
\end{enumerate}

\item Какую сходимость определяет в пространстве $C[a, b]$ задание топологии с помощью фундаментальной системы окрестностей вида:
\begin{enumerate}[label=\alph*)]
\item $\mathcal{U}_{\epsilon}(x_{0}) = \{x \in C[a, b] : |x(t) - x_{0}(t)| < \epsilon \ \forall t \in [a, b]\}$, \\ где $x_{0}(t) \in C[a, b]$ — фиксированная функция;
\item $\mathcal{U}_{\epsilon, F}(x_{0}) = \{x \in C[a, b] : |x(t) - x_{0}(t)| < \epsilon \ \forall t \in F\}$, \\ где $x_{0}(t) \in C[a, b]$ — фиксированная функция, а $F$ — фиксированное конечное множество из $[a, b]$.
\end{enumerate}

\item Какая из топологий, определенных в предыдущем номере, сильнее?

\item На пространстве функций, определенных на $[a, b]$, ввести фундаментальные системы окрестностей, определяющие сходимости:
\begin{enumerate}[label=\alph*)]
\item по мере;
\item почти всюду.
\end{enumerate}

\item Построить слабейшую из топологий на $\mathbb{R}$, содержащих все множества $\mathcal{B}$ вида:
\begin{enumerate}[label=\alph*)]
\item $\mathcal{B} = \{(a, b) : a, b \in \mathbb{R}\}$;
\item $\mathcal{B} = \{[a, b] : a, b \in \mathbb{R}\}$;
\item $\mathcal{B} = \{[a, b) : a, b \in \mathbb{R}\}$;
\item $\mathcal{B} = \{(a, b] : a, b \in \mathbb{R}\}$;
\item $\mathcal{B} = \{(a, +\infty) : a \in \mathbb{R}\}$;
\item $\mathcal{B} = \{[a, +\infty) : a \in \mathbb{R}\}$;
\item $\mathcal{B} = \{(-\infty, a) : a \in \mathbb{R}\}$;
\item $\mathcal{B} = \{(a, b) : a, b \in \mathbb{R}, a > 0, b > 0\}$;
\item $\mathcal{B} = \{[n, m] : n, m \in \mathbb{Z}\}$.
\end{enumerate}

\item Найти пересечения топологий, построенных в задачах 9a и 9b.
\item Найти пересечения топологий, построенных в задачах 9c и 9d.
\item Сравнить топологии, построенных в задачах 9e и 9f.

\item Пусть $\{\tau_{n}\}_{n=1}^{\infty}$ — последовательность топологий на $X$, причем $\tau_{1} \subset \tau_{2} \subset \dots$ Будет ли топологией $\bigcup_{n=1}^{\infty} \tau_{n}$ на $X$?

\item Сходится ли последовательность $\{\frac{1}{n}\}_{n=1}^{\infty}$ в топологиях, построенных в задачах 9a, 9b, 9c, 9d, 9e?

\item Сходится ли последовательность $\{\frac{(-1)^{n}}{n}\}_{n=1}^{\infty}$ в топологиях, построенных в задачах 9a, 9b, 9c, 9d?

\item Сходится ли последовательность $\{(-1)^{n}\}_{n=1}^{\infty}$ в топологиях, построенных в задачах 9a, 9b, 9c, 9d, 9e?

\item Сходится ли последовательность $\{n \cdot (-1)^{n}\}_{n=1}^{\infty}$ в топологиях, построенных в задачах 9a, 9b, 9c, 9d?

\item Описать сходящиеся последовательности топологического пространства $(X, \tau)$, если:
\begin{enumerate}[label=\alph*)]
\item $\tau$ — слабейшая топология;
\item $\tau$ — сильнейшая топология.
\end{enumerate}

\item Описать сходящиеся последовательности топологического пространства $\mathbb{R}$ с каждой из топологий, построенных в задачах 9a, 9b, 9c, 9d.
\end{enumerate}

\section{МЕТРИЧЕСКИЕ ПРОСТРАНСТВА}

\subsection*{КОНТРОЛЬНЫЕ ВОПРОСЫ}
\begin{enumerate}
\item Дать определение метрического пространства.
\item Как определяется метрика в пространствах $\mathbb{R}^{n}$, $C[a, b]$, $C^{(m)}[a, b]$, $l^{p}$, $L^{p}(a, b)$, $L^{\infty}(a, b)$, $l^{\infty}=m$, $s$?
\item Дать определение сходящейся последовательности элементов метрического пространства.
\item Дать определения открытого и замкнутого множеств и замыкания множества.
\end{enumerate}

\subsection*{ЗАДАЧИ}
\begin{enumerate}
\item Какие из нижеприведенных формул определяют метрику в $X$:
\begin{enumerate}[label=\alph*)]
\item $\rho(x, y) = \max_{0\le t\le1} |x(t) - y(t)|$, $X = C[0, 2]$;
\item $\rho(x, y) = \cos^{2}(x - y)$, $X = \mathbb{R}^{1}$;
\item $\rho(x, y) = |\operatorname{arctg} x - \operatorname{arctg} y|$, $X = \mathbb{R}^{1}$;
\item $\rho(x, y) = |\sin(x - y)|$, $X = (-\frac{\pi}{2}, \frac{\pi}{2})$;
\item $\rho(x, y) = \sum_{k=1}^{n} |\xi_{k} - \eta_{k}|$, $X = \mathbb{R}^{n}$, \\ где $x = (\xi_{1}, \dots, \xi_{n}) \in \mathbb{R}^{n}$, $y = (\eta_{1}, \dots, \eta_{n}) \in \mathbb{R}^{n}$;
\item $\rho(x, y) = \sum_{k=1}^{n} |\xi_{k} - \eta_{k}|^{2}$, $X = \mathbb{R}^{n}$;
\item $\rho(x, y) = \sum_{k=1}^{\infty} |\xi_{k} - \eta_{k}|$, $X = l^{\infty}$;
\item $\rho(x, y) = \operatorname{arctg}|x - y|$, $X = \mathbb{R}^{1}$;
\item $\rho(x, y) = \left( \sum_{k=1}^{\infty} |\xi_{k} - \eta_{k}|^{p} \right)^{1/p}$, $X = l^{p}$, $0 < p < 1$;
\item $\rho(x, y) = \left( \int_{a}^{b} |x(t) - y(t)|^{p} \, dt \right)^{1/p}$, $X = L^{p}(a, b)$, $0 < p < 1$;
\item $\rho(x, y) = \max_{0\le t\le2} |x^{\prime}(t) - y^{\prime}(t)|$, $X = C^{(1)}[0, 2]$?
\end{enumerate}

\item Выяснить, сходится ли в метрическом пространстве $X$ последовательность $\{x_{n}\}_{n=1}^{\infty}$. Найти предел, если сходится.
\begin{enumerate}[label=\alph*)]
\item $X = C[0, 1]$, $x_{n}(t) = t^{n}$;
\item $X = C[0, 1]$, $x_{n}(t) = t^{n} - t^{n+1}$;
\item $X = l^{2}$, $x^{(n)} = (\frac{1}{n}, 0, 0, \dots, 0, 1, 0, 0, \dots)$ (1 на $n$-ом месте);
\item $X = l^{2}$, $x^{(n)} = (1, 0, \dots, 0, \frac{1}{n}, 0, \dots)$ ($\frac{1}{n}$ на $n$-ом месте);
\item $X = L^{2}(0, 1)$, $x_{n}(t) = t^{n} - t^{2n}$;
\item $X = L^{2}(0, 1)$, $x_{n}(t) =
\begin{cases} 1-nt, & \text{если } t\in[0, \frac{1}{n}], \\ 0, & \text{если } t\in(\frac{1}{n}, 1];
\end{cases}$
\item $X = C[0, 2]$, $x_{n}(t) = e^{-t/n}$;
\item $X = C^{(1)}[0, 1]$, $x_{n}(t) = \frac{t^{n+1}}{n+1} - \frac{t^{n+2}}{n+2}$;
\item $X = L^{1}(0, 3)$, $x_{n}(t) = n e^{-nt}$;
\item $X = s$, $x^{(n)} = (1, 2, 3, \dots, n, 0, 0, \dots)$.
\end{enumerate}

\item В каком из пространств $l^{p}$, $l^{\infty}$, $s$ и при каком фиксированном $\alpha$ сходится последовательность
$$x^{(n)} = \left(\frac{1}{n^{\alpha}}, \frac{1}{n^{\alpha}}, \dots, \frac{1}{n^{\alpha}}, 0, 0, \dots \right)?$$ (где $\frac{1}{n^{\alpha}}$ повторяется $n$ раз)

\newpage
\item Доказать, что для любых четырех точек $x, y, z, t$ метрического пространства справедливо „неравенство четырёхугольника“:
$$|\rho(x, z) - \rho(y, t)| \le \rho(x, y) + \rho(z, t).$$


\begin{center}\small\textbf{Формализация на Lean 4}\end{center}
\inputminted{lean}{FunAn/Section07_MetricSpaces/Task04.lean}

\newpage
\item Доказать, что $\rho(x, A) = 0$ тогда и только тогда, когда $x \in \overline{A}$, где $A$ — подмножество метрического пространства $X$.


\begin{center}\small\textbf{Формализация на Lean 4}\end{center}
\inputminted{lean}{FunAn/Section07_MetricSpaces/Task05.lean}

\newpage
\item Доказать непрерывность функции $f(x) = \rho(x, A)$.


\begin{center}\small\textbf{Формализация на Lean 4}\end{center}
\inputminted{lean}{FunAn/Section07_MetricSpaces/Task06.lean}

\newpage
\item Доказать, что множество $\mathcal{M} = \{x : \rho(x, A) < \epsilon\}$ открыто, а множество $\mathcal{N} = \{x : \rho(x, A) \le \epsilon\}$ замкнуто.\\
\textit{Указание.} Доказать, что $\mathcal{M} = \bigcup_{y\in A} B(y, \epsilon)$. Использовать задачу 6.


\begin{center}\small\textbf{Формализация на Lean 4}\end{center}
\inputminted{lean}{FunAn/Section07_MetricSpaces/Task07.lean}

\item Приведите пример последовательности дважды непрерывно дифференцируемых на отрезке $[a, b]$ функций, которая сходится в пространстве $C[a, b]$, но не сходится в пространстве $C^{(2)}[a, b]$.
\end{enumerate}

\section{ПОЛНЫЕ МЕТРИЧЕСКИЕ ПРОСТРАНСТВА. ПОПОЛНЕНИЕ}

\subsection*{КОНТРОЛЬНЫЕ ВОПРОСЫ}
\begin{enumerate}
\item Дать определения фундаментальной последовательности и полного метрического пространства.
\item Привести примеры полных и неполных метрических пространств.
\item Сформулировать принцип вложенных шаров.
\item Дать определение пополнения метрического пространства.
\end{enumerate}

\subsection*{ЗАДАЧИ}
\begin{enumerate}
\newpage
\item На множестве $\mathbb{R}^{1}$ введены метрики
$$\rho(x, y) = |x - y|, \quad \rho_{1}(x, y) = |\operatorname{arctg} x - \operatorname{arctg} y|.$$
Показать, что в пространствах $(\mathbb{R}^{1}, \rho)$ и $(\mathbb{R}^{1}, \rho_{1})$ сходящиеся последовательности одни и те же, а фундаментальные различны.


\begin{center}\small\textbf{Формализация на Lean 4}\end{center}
\inputminted{lean}{FunAn/Section08_CompleteMetric/Task01.lean}

\item Доказать полноту следующих метрических пространств:
\begin{enumerate}[label=\alph*)]
\item $l_{\{k\}}^{p} = \left\{x = (\xi_{1}, \xi_{2}, \dots) : \xi_{k} \in \mathbb{C}, \sum_{k=1}^{\infty} k|\xi_{k}|^{p} < \infty\right\}$, $1 \le p < \infty$; \\
$\rho(x, y) = \left( \sum_{k=1}^{\infty} k|\xi_{k} - \eta_{k}|^{p} \right)^{1/p}$;
\item $C^{(m)}[a, b]$, $\rho(x, y) = \sum_{k=0}^{m} \max_{t\in[a, b]} |x^{(k)}(t) - y^{(k)}(t)|$;
\item $C^{\infty}[a, b]$, $\rho(x, y) = \sum_{k=0}^{\infty} \frac{1}{2^{k}} \cdot \frac{\max_{t\in[a, b]} |x^{(k)}(t) - y^{(k)}(t)|}{1 + \max_{t\in[a, b]} |x^{(k)}(t) - y^{(k)}(t)|}$;
\item $C(\mathbb{R}^{1})$, $\rho(x, y) = \sum_{k=1}^{\infty} \frac{1}{2^{k}} \cdot \frac{\max_{|t|\le k} |x(t) - y(t)|}{1 + \max_{|t|\le k} |x(t) - y(t)|}$;
\item $l_{\infty} = \left\{x = (\xi_{1}, \xi_{2}, \dots) : \sup_{n\ge1} \left| \sum_{k=1}^{n} \xi_{k} \right| < \infty \right\}$, \\ $\rho(x, y) = \sup_{n\ge1} \left| \sum_{k=1}^{n} (\xi_{k} - \eta_{k}) \right|$;
\item $c_{0} = \left\{x = (\xi_{1}, \xi_{2}, \dots) : \xi_{k} \in \mathbb{C}, \xi_{k} \rightarrow 0 \text{ при } k \rightarrow \infty\right\}$, \\ $\rho(x, y) = \max_{k\ge1} |\xi_{k} - \eta_{k}|$;
\item $X = \{x \in C(\mathbb{R}^{1}) : x(t) \rightarrow 0 \text{ при } |t| \rightarrow \infty\}$, \\ $\rho(x, y) = \max_{t\in\mathbb{R}} |x(t) - y(t)|$;
\newpage
\item $X = \{x \in C[0, 1] : x(0) = 0\}$, $\rho(x, y) = \max_{0\le t\le1} |x(t) - y(t)|$.
\end{enumerate}


\begin{center}\small\textbf{Формализация на Lean 4}\end{center}
\inputminted{lean}{FunAn/Section08_CompleteMetric/Task02.lean}

\newpage
\item Доказать полноту пространства $\mathbb{N}$ с метрикой
$$\rho(n, m) =
\begin{cases} 0, & \text{если } n=m, \\ 1 + \frac{1}{n+m}, & \text{если } n \ne m.
\end{cases}$$


\begin{center}\small\textbf{Формализация на Lean 4}\end{center}
\inputminted{lean}{FunAn/Section08_CompleteMetric/Task03.lean}

\newpage
\item Доказать, что шары $B[n, 1+\frac{1}{2n}]$ пространства $\mathbb{N}$ (см. задачу 3) замкнуты и вложены, но не имеют общей точки. Почему это не противоречит принципу вложенных шаров?


\begin{center}\small\textbf{Формализация на Lean 4}\end{center}
\inputminted{lean}{FunAn/Section08_CompleteMetric/Task04.lean}

\item Привести пример метрического пространства, в котором существует последовательность замкнутых вложенных шаров, не имеющих общей точки, радиусы которых стремятся к нулю.

\item Доказать, что следующие метрические пространства не являются полными, построить их пополнения.
\begin{enumerate}[label=\alph*)]
\item $C[0, 1]$, $\rho(x, y) = \int_{0}^{1} |x(t) - y(t)| \, dt$;
\item $X = \{x = (\xi_{1}, \xi_{2}, \dots) : \xi_{k} \in \mathbb{C}, \xi_{k}=0 \text{ при } k > k_{0}(x)\}$, \\ $\rho(x, y) = \max_{k\ge1} |\xi_{k} - \eta_{k}|$;
\item $C_{0}(\mathbb{R}^{1}) = \{x(t) \in C(\mathbb{R}^{1}) : x(t) \equiv 0 \text{ при } |t| > t_{0}(x)\}$, \\ $\rho(x, y) = \max_{t\in\mathbb{R}} |x(t) - y(t)|$;
\item $\mathbb{R}^{1}$, $\rho(x, y) = |\operatorname{arctg} x - \operatorname{arctg} y|$;
\item $\mathbb{Z}$, $\rho(p, q) = |e^{ip} - e^{iq}|$;
\item $X = \{x \in C^{\infty}[0, 1] : x^{(k)}(0) = x^{(k)}(1) = 0 \ \forall k \ge 0\}$, \\ $\rho(x, y) = \max_{0\le t\le1} |x(t) - y(t)|$;
\newpage
\item $C[0, 1]$, $\rho(x, y) = \max_{0\le t\le1} t|x(t) - y(t)|$.
\end{enumerate}
\textit{Указание.} Рассмотреть
$$x_{n}(t) =
\begin{cases} n, & \text{если } 0 \le t \le \frac{1}{n^{2}}, \\ \frac{1}{\sqrt{t}}, & \text{если } t > \frac{1}{n^{2}}.
\end{cases}$$
См. также задачу 8.

\begin{center}\small\textbf{Формализация на Lean 4}\end{center}
\inputminted{lean}{FunAn/Section08_CompleteMetric/Task06.lean}

\end{enumerate}

\section{СЕПАРАБЕЛЬНЫЕ МЕТРИЧЕСКИЕ ПРОСТРАНСТВА}

\subsection*{КОНТРОЛЬНЫЕ ВОПРОСЫ}
\begin{enumerate}
\item Дать определение счетного множества, множества мощности континуума. Приведите примеры.
\item Определите плотные, всюду плотные, нигде не плотные множества.
\item Дать определение сепарабельного метрического пространства. Привести примеры.
\end{enumerate}

\subsection*{ЗАДАЧИ}
\begin{enumerate}
\newpage
\item Доказать счетность множества финитных числовых последовательностей с рациональными членами.


\begin{center}\small\textbf{Формализация на Lean 4}\end{center}
\inputminted{lean}{FunAn/Section09_Separable/Task01.lean}

\newpage
\item Доказать счетность множества многочленов с рациональными коэффициентами.


\begin{center}\small\textbf{Формализация на Lean 4}\end{center}
\inputminted{lean}{FunAn/Section09_Separable/Task02.lean}

\newpage
\item Доказать, что множество $A=\{z=e^{in} : n\in\mathbb{Z}\}$ плотно в множестве $E=\{z\in\mathbb{C} : |z|=1\}$.


\begin{center}\small\textbf{Формализация на Lean 4}\end{center}
\inputminted{lean}{FunAn/Section09_Separable/Task03.lean}

\newpage
\item Доказать, что множество
$$B=\{x=(\xi_{1},\xi_{2},\dots,\xi_{n},0,0,\dots) : \xi_{k}\in\mathbb{C}\}$$
нигде не плотно в $l^{p}$ ($n$ — фиксировано).


\begin{center}\small\textbf{Формализация на Lean 4}\end{center}
\inputminted{lean}{FunAn/Section09_Separable/Task04.lean}

\newpage
\item Доказать, что пространства $\mathbb{R}^{n}$, $C[a,b]$, $l^{p}$, $s$, $C^{(m)}[a,b]$, $L^{p}(a,b)$ сепарабельны $(1\le p<\infty)$.


\begin{center}\small\textbf{Формализация на Lean 4}\end{center}
\inputminted{lean}{FunAn/Section09_Separable/Task05.lean}

\newpage
\item Доказать, что $l^{\infty}$ не является сепарабельным.


\begin{center}\small\textbf{Формализация на Lean 4}\end{center}
\inputminted{lean}{FunAn/Section09_Separable/Task06.lean}

\newpage
\item Доказать, что пространство $M[a,b]$ ограниченных на $[a, b]$ функций с метрикой $\rho(x,y)=\sup_{a\le t\le b}|x(t)-y(t)|$ не является сепарабельным.


\begin{center}\small\textbf{Формализация на Lean 4}\end{center}
\inputminted{lean}{FunAn/Section09_Separable/Task07.lean}

\newpage
\item Доказать, что совокупность подпоследовательностей последовательности натуральных чисел несчетно.


\begin{center}\small\textbf{Формализация на Lean 4}\end{center}
\inputminted{lean}{FunAn/Section09_Separable/Task08.lean}

\newpage
\item Доказать, что мощность сепарабельного пространства не превосходит мощности континуума.\\
\textit{Указание.} Использовать задачу 8.


\begin{center}\small\textbf{Формализация на Lean 4}\end{center}
\inputminted{lean}{FunAn/Section09_Separable/Task09.lean}

\newpage
\item Доказать сепарабельность пространства $l_{\{k\}}^{p}$ (см. задачу 2a).


\begin{center}\small\textbf{Формализация на Lean 4}\end{center}
\inputminted{lean}{FunAn/Section09_Separable/Task10.lean}

\newpage
\item Доказать несепарабельность пространства $l_{\infty}$ (см. задачу 2e).

\begin{center}\small\textbf{Формализация на Lean 4}\end{center}
\inputminted{lean}{FunAn/Section09_Separable/Task11.lean}

\end{enumerate}

\section{ПРИНЦИП СЖИМАЮЩИХ ОТОБРАЖЕНИЙ}

\subsection*{КОНТРОЛЬНЫЕ ВОПРОСЫ}
\begin{enumerate}
\item Дать определения непрерывного оператора, сжимающего оператора, неподвижной точки оператора.
\item Сформулировать принцип сжимающих отображений и его обобщение.
\item Какое уравнение называется интегральным уравнением Вольтерра? Фредгольма?
\end{enumerate}

\subsection*{ЗАДАЧИ}
\begin{enumerate}
\item Являются ли непрерывными операторы:
\begin{enumerate}[label=\alph*)]
\item $(Ax)(t)=\int_{0}^{t}x^{2}(s) \, ds$, $A:C[0,1]\rightarrow C[0,1]$;
\item $(Ax)(t)=x^{\prime}(t)$, $A:C^{(1)}[0,2]\rightarrow C[0,2]$;
\item $(Ax)(t)=x^{\prime}(t)$, $A:C^{(1)}[0,2]\rightarrow C[0,2]$, \\ причём расстояние в $C^{(1)}[0,2]$ определено так же, как в $C[0,2]$.
\end{enumerate}

\newpage
\item Проверьте, что функция $f(x)=\sqrt{1+x^{2}}$, определённая на $[0, +\infty)$, удовлетворяет условию $|f(x)-f(y)|<|x-y|$ при $x\ne y$, однако уравнение $x=f(x)$ не имеет решений.


\begin{center}\small\textbf{Формализация на Lean 4}\end{center}
\inputminted{lean}{FunAn/Section10_Contraction/Task02.lean}

\newpage
\item Покажите, что отображение $f(x)=\frac{x^{2}+2}{2x}$ является сжимающим на отрезке $[1, 2]$. Вычислите несколько итераций по формуле $x_{0}=1$, $x_{k+1}=f(x_{k})$, $k=0,1,\dots$, и сравните их с табличным значением $\sqrt{2}$ (решение уравнения $x=f(x)$, лежащее на отрезке $[1,2]$).


\begin{center}\small\textbf{Формализация на Lean 4}\end{center}
\inputminted{lean}{FunAn/Section10_Contraction/Task03.lean}

\newpage
\item Пусть $f(x)=\frac{\pi}{2}+x-\operatorname{arctg} x$, $x\in\mathbb{R}^{1}$. Доказать, что для любых $x, y\in\mathbb{R}^{1}$ найдется постоянная $\alpha<1$ такая, что
$$|f(x)-f(y)|\le\alpha|x-y|,$$
но $f$ не имеет неподвижных точек. Почему это не противоречит принципу сжимающих отображений?


\begin{center}\small\textbf{Формализация на Lean 4}\end{center}
\inputminted{lean}{FunAn/Section10_Contraction/Task04.lean}

\item При каких $\lambda$ и $\beta$ отображение $(Ax)(t)=\lambda\cdot x(t^{\beta})$ является сжимающим в $C[0,1]$? В $L^{2}(0,1)$?

\newpage
\item Доказать, что если $\sup_{j\ge1}\sum_{i=1}^{\infty}|a_{ij}|<1$, то бесконечная система линейных алгебраических уравнений
$$x_{i}=\sum_{j=1}^{\infty}a_{ij}x_{j}+b_{i} \quad (i=1,2,\dots)$$
имеет единственное решение $x=(x_{1},x_{2},\dots)\in l^{1}$ для любой последовательности $b=(b_{1},b_{2},\dots)\in l^{1}$.


\begin{center}\small\textbf{Формализация на Lean 4}\end{center}
\inputminted{lean}{FunAn/Section10_Contraction/Task06.lean}

\newpage
\item Доказать, что если $\sup_{i\ge1}\sum_{j=1}^{\infty}|a_{ij}|<1$, то утверждение задачи 6 верно в пространстве $l^{\infty}$.


\begin{center}\small\textbf{Формализация на Lean 4}\end{center}
\inputminted{lean}{FunAn/Section10_Contraction/Task07.lean}

\newpage
\item Проверить, что интегральное уравнение Вольтерра
$$x(t)=\int_{0}^{t}K(t,s)x(s) \, ds,$$
где
$$K(t,s)=
\begin{cases} s\cdot e^{-1+1/t^{2}}, & \text{если } 0\le s\le t\cdot e^{1-1/t^{2}}, \\ t, & \text{если } t\cdot e^{1-1/t^{2}}\le s\le t\le1,
\end{cases}$$
имеет бесконечно много решений $x(t)=\frac{c}{t}$. Почему это не противоречит теореме о существовании и единственности решения уравнения Вольтерра?


\begin{center}\small\textbf{Формализация на Lean 4}\end{center}
\inputminted{lean}{FunAn/Section10_Contraction/Task08.lean}

\item При каких $\lambda$ применим принцип сжимающих отображений в пространствах $L^{2}(a,b)$ и $C[a,b]$ к уравнению Фредгольма 2-го рода
$$x(t)=\lambda\int_{a}^{b}K(t,s)x(s) \, ds+y(t) ?$$
Найти точное решение уравнения. Для нахождения приближенного решения составить программу. В программе предусмотреть:
\begin{enumerate}[label=\alph*)]
\item табулирование точного решения с шагом $h=\frac{b-a}{2n}$;
\item вычисление номера $L$ итерации, обеспечивающей точность $0,01$ в метрике $C[a,b]$, когда за первое приближение берется $x_{1}(t)\equiv y(t)$. Заметим, что $L$ определяется из неравенства
$$\frac{\theta^{L-1}}{1-\theta}\cdot\rho(x_{1},x_{2})<0,01,$$
где $\theta$ — коэффициент сжатия;
\item сравнение точного решения с приближенным, т.е. вычисление величины $R=\max_{0\le k\le2n}|x(t_{k})-\overline{x}(t_{k})|$, где $x(t_{k})$ — значения точного решения, $\overline{x}(t_{k})$ — значения приближенного решения.
\end{enumerate}
Для вычисления $\overline{x}(t_{k})$ разбиваем промежуток $[a,b]$ на $2n$ частей с помощью точек $t_{k}=a+kh$. Тогда
$$x(t_{k})=\lambda\int_{a}^{b}K(t_{k},s)x(s) \, ds+y(t_{k}), \quad k=0,1,\dots,2n.$$
С помощью формулы Симпсона заменим интегралы на приближенные значения. Получим:
\begin{align*}
x(t_{k}) &= \frac{\lambda h}{3} \Bigg[ K(t_{k},s_{0})x(s_{0}) + K(t_{k},s_{2n})x(s_{2n}) \\
&\quad + 4\sum_{m=1}^{n}K(t_{k},s_{2m-1})x(s_{2m-1}) + 2\sum_{m=1}^{n-1}K(t_{k},s_{2m})x(s_{2m}) \Bigg] + y(t_{k}),
\end{align*}
$k=0,1,\dots,2n$. Полученная система линейных алгебраических уравнений решается методом итераций.

Задания приведены в таблице 2. В отчете показать совокупность значений $\lambda$, допускающих применение принципа сжатых отображений в пространствах $L^{2}[a,b]$, $C[a,b]$, точное решение уравнения, значения чисел $L$ и $R$, а также привести программу и заполненную таблицу 1.
\end{enumerate}

\begin{table}[ht!]
\centering
\caption*{Таблица 1.}
\begin{tabular}{|c|c|c|c|c|}
\hline
$k$ & $0$ & $1$ & $2$ & $\dots \quad 2n$ \\
\hline
$x(t_{k})$ & & & & \\
\hline
$\overline{x}(t_{k})$ & & & & \\
\hline
\end{tabular}
\end{table}

\begin{table}[ht!]
\centering
\caption*{Таблица 2.}
\begin{tabular}{|c|c|c|c|c|c|c|}
\hline
№ & $K(t,s)$ & $y(t)$ & $a$ & $b$ & $\lambda$ & $n$ \\
\hline
1. & $t^{2}s$ & $\cos 3t$ & $0$ & $1$ & $1$ & $8$ \\
\hline
2. & $t^{2}s$ & $\cos 3t$ & $0$ & $1$ & $1$ & $10$ \\
\hline
3. & $ts$ & $e^{t}$ & $-1$ & $1$ & $0,5$ & $10$ \\
\hline
4. & $ts$ & $e^{t}$ & $0,5$ & $1,5$ & $0,5$ & $10$ \\
\hline
5. & $t^{2}s^{2}$ & $\frac{\sin t}{t}$ & $0,3$ & $1,2$ & $0,6$ & $9$ \\
\hline
6. & $e^{t-s}$ & $1$ & $0,1$ & $0,9$ & $0,5$ & $8$ \\
\hline
7. & $e^{t-s}$ & $t$ & $0,1$ & $0,8$ & $0,5$ & $7$ \\
\hline
8. & $\cos \pi(t-s)$ & $1$ & $0$ & $1$ & $0,5$ & $10$ \\
\hline
9. & $\cos t \sin s$ & $1$ & $0$ & $2$ & $0,25$ & $10$ \\
\hline
10. & $\sin t \cos s$ & $1$ & $-0,5$ & $0,5$ & $0,25$ & $5$ \\
\hline
11. & $e^{2t}\cos s$ & $\sin t$ & $0$ & $0,5$ & $0,6$ & $5$ \\
\hline
12. & $e^{2t}\sin s$ & $\cos t$ & $0$ & $0,5$ & $0,4$ & $10$ \\
\hline
13. & $e^{t}\cos s$ & $2\sin t$ & $0$ & $0,9$ & $0,5$ & $8$ \\
\hline
14. & $e^{t}\sin s$ & $2\cos t$ & $0$ & $0,8$ & $0,5$ & $8$ \\
\hline
15. & $t\sin s$ & $e^{t}$ & $0,2$ & $0,8$ & $0,4$ & $6$ \\
\hline
16. & $t\cos s$ & $e^{-t}$ & $0,1$ & $0,9$ & $0,3$ & $8$ \\
\hline
17. & $t^{2}e^{s}$ & $t$ & $-2$ & $-1$ & $2$ & $6$ \\
\hline
\end{tabular}
\end{table}

\section{НОРМИРОВАННЫЕ ПРОСТРАНСТВА}

\subsection*{КОНТРОЛЬНЫЕ ВОПРОСЫ}
\begin{enumerate}
\item Дать определения линейного пространства, нормы, нормированного пространства, банахова пространства.
\item Дать определение эквивалентных норм.
\item Какая система элементов линейного пространства называется линейно независимой? Дать определение размерности линейного пространства. Привести примеры.
\end{enumerate}

\subsection*{ЗАДАЧИ}
\begin{enumerate}
\newpage
\item Доказать, что аксиома треугольника в определении нормы эквивалентна выпуклости единичного шара.

\begin{center}\small\textbf{Формализация на Lean 4}\end{center}
\inputminted{lean}{FunAn/Section11_NormedSpaces/Task01.lean}

\item Можно ли в пространстве $s$ ввести норму, определяющую известную метрику?
\newpage
\item Доказать, что нормированное пространство $X$ является банаховым пространством тогда и только тогда, когда всякий ряд $\sum_{k=1}^{\infty}x_{k}$, для которого $\sum_{k=1}^{\infty}||x_{k}||<\infty$, сходится в $X$.

\begin{center}\small\textbf{Формализация на Lean 4}\end{center}
\inputminted{lean}{FunAn/Section11_NormedSpaces/Task03.lean}

\newpage
\item Доказать, что $l^{p}\subset l^{q}$ при $p<q$.

\begin{center}\small\textbf{Формализация на Lean 4}\end{center}
\inputminted{lean}{FunAn/Section11_NormedSpaces/Task04.lean}

\newpage
\item Доказать, что $L^{p}(a,b)\supset L^{q}(a,b)$ при $p<q$.

\begin{center}\small\textbf{Формализация на Lean 4}\end{center}
\inputminted{lean}{FunAn/Section11_NormedSpaces/Task05.lean}

\item Можно ли в $C^{(1)}[a,b]$ ввести нормы по формулам:
\begin{enumerate}[label=\alph*)]
\item $||x||=\max_{t\in[a,b]}|x^{\prime}(t)|$;
\item $||x||=|x(b)-x(a)|+\max_{t\in[a,b]}|x^{\prime}(t)|$;
\item $||x||=|x(a)|+\max_{a\le t\le b}|x^{\prime}(t)|$;
\item $||x||=\int_{a}^{b}|x(t)| \, dt+\max_{a\le t\le b}|x^{\prime}(t)|$.
\end{enumerate}
\newpage
\item Доказать, что нормы $||x||=\max_{a\le t\le b}|x(t)|$ и $||x||_{1}=\int_{a}^{b}|x(t)| \, dt$ не эквивалентны на $C[a,b]$.

\begin{center}\small\textbf{Формализация на Lean 4}\end{center}
\inputminted{lean}{FunAn/Section11_NormedSpaces/Task07.lean}

\newpage
\item Доказать, что $C[a,b]$ с нормой $||x||=\int_{a}^{b}|x(t)| \, dt$ не полно.

\begin{center}\small\textbf{Формализация на Lean 4}\end{center}
\inputminted{lean}{FunAn/Section11_NormedSpaces/Task08.lean}

\newpage
\item Показать, что нормы $\max(||x_{1}||,||x_{2}||)$, $||x_{1}||+||x_{2}||$, $\sqrt{||x_{1}||^{2}+||x_{2}||^{2}}$ эквивалентны на $E\times E$, где $|| \cdot ||$ — норма в $E$.

\begin{center}\small\textbf{Формализация на Lean 4}\end{center}
\inputminted{lean}{FunAn/Section11_NormedSpaces/Task09.lean}

\newpage
\item Доказать, что конечномерное подпространство нормированного пространства всегда замкнуто.

\begin{center}\small\textbf{Формализация на Lean 4}\end{center}
\inputminted{lean}{FunAn/Section11_NormedSpaces/Task10.lean}

\newpage
\item Если $L$ — подпространство нормированного пространства $E$, то для любого $\alpha<1$ существует $y\in E$, $||y||=1$, такой, что $\rho(y,L)>\alpha$. Доказать.

\begin{center}\small\textbf{Формализация на Lean 4}\end{center}
\inputminted{lean}{FunAn/Section11_NormedSpaces/Task11.lean}

\newpage
\item В бесконечномерном нормированном пространстве существует бесконечное ограниченное множество, элементы которого удалены друг от друга на расстояние $>0,5$. Доказать.\\
\textit{Указание.} Использовать задачу 11.

\begin{center}\small\textbf{Формализация на Lean 4}\end{center}
\inputminted{lean}{FunAn/Section11_NormedSpaces/Task12.lean}

\newpage
\item Показать на примере $E=C[-1,1]$,
$$L=\left\{x(t)\in C[-1,1] : \int_{-1}^{0}x(t) \, dt=\int_{0}^{1}x(t) \, dt\right\},$$
что постоянную $\alpha$ в задаче 11 нельзя заменить на 1.

\begin{center}\small\textbf{Формализация на Lean 4}\end{center}
\inputminted{lean}{FunAn/Section11_NormedSpaces/Task13.lean}

\item В пространстве $\mathbb{R}^{2}$ с нормой $||x||=\max\{|\xi_{1}|,|\xi_{2}|\}$ найти подпространство $L$ и точку $x_{0}$ такие, что ближайший элемент в $L$ до $x_{0}$ не единственный.
\newpage
\item Доказать, что множество $\mathcal{M}=\{x\in l^{1}:||x||_{2}\le1\}$ не является замкнутым множеством в $l^{2}$, если $||x||_{2}=\left(\sum_{k=1}^{\infty}|\xi_{k}|^{2}\right)^{1/2}$.

\begin{center}\small\textbf{Формализация на Lean 4}\end{center}
\inputminted{lean}{FunAn/Section11_NormedSpaces/Task15.lean}

\newpage
\item Доказать, что множество $\mathcal{M}=\{x\in C^{(1)}[-1,1]:||x||_{1}\le1\}$ не является замкнутым в $C[-1,1]$ с $||x||=\max_{|t|\le1}|x(t)|$, где $||x||_{1}=\max_{|t|\le1}|x(t)|+\max_{|t|\le1}|x^{\prime}(t)|$.\\
\textit{Указание.} Рассмотреть $x_{n}(t)=\frac{1}{3}\cdot\sqrt{t^{2}+\frac{1}{n^{2}}}$.

\begin{center}\small\textbf{Формализация на Lean 4}\end{center}
\inputminted{lean}{FunAn/Section11_NormedSpaces/Task16.lean}

\newpage
\item Доказать, что замкнутый единичный шар пространства $l^{1}$ — замкнутое множество пространства $l^{2}$.

\begin{center}\small\textbf{Формализация на Lean 4}\end{center}
\inputminted{lean}{FunAn/Section11_NormedSpaces/Task17.lean}

\item Является ли пространство $c$ всех сходящихся к конечному пределу числовых последовательностей банаховым пространством, если $||x||=\sup_{k\ge1}|\xi_{k}|$?
\end{enumerate}

\section{ГИЛЬБЕРТОВЫ ПРОСТРАНСТВА}

\subsection*{КОНТРОЛЬНЫЕ ВОПРОСЫ}
\begin{enumerate}
\item Дать определение гильбертова пространства.
\item Дать определение ортогонального дополнения множества.
\item Сформулировать теорему о проекциях.
\item Какая система элементов называется ортогональной, ортонормированной, полной, замкнутой?
\end{enumerate}

\subsection*{ЗАДАЧИ}
\begin{enumerate}
\newpage
\item Доказать, что формула $(x,y)=\sum_{k=1}^{\infty}\alpha_{k}\xi_{k}\overline{\eta}_{k}$, $0<\alpha_{k}\le1$, определяет скалярное произведение в $l^{2}$.

\begin{center}\small\textbf{Формализация на Lean 4}\end{center}
\inputminted{lean}{FunAn/Section12_HilbertSpaces/Task01.lean}

\newpage
\item Доказать равенство параллелограмма:
$$||x+y||^{2}+||x-y||^{2}=2||x||^{2}+2||y||^{2}, \quad \forall x,y\in H.$$

\begin{center}\small\textbf{Формализация на Lean 4}\end{center}
\inputminted{lean}{FunAn/Section12_HilbertSpaces/Task02.lean}

\newpage
\item Доказать, что пространства $l^{p}$ и $L^{p}(a,b)$ являются гильбертовыми тогда и только тогда, когда $p=2$.

\begin{center}\small\textbf{Формализация на Lean 4}\end{center}
\inputminted{lean}{FunAn/Section12_HilbertSpaces/Task03.lean}

\newpage
\item Доказать, что пространство $C[0,\pi]$ не является гильбертовым.

\begin{center}\small\textbf{Формализация на Lean 4}\end{center}
\inputminted{lean}{FunAn/Section12_HilbertSpaces/Task04.lean}

\newpage
\item Доказать поляризационное тождество:
$$4(x,y)=||x+y||^{2}-||x-y||^{2}+i||x+iy||^{2}-i||x-iy||^{2}.$$

\begin{center}\small\textbf{Формализация на Lean 4}\end{center}
\inputminted{lean}{FunAn/Section12_HilbertSpaces/Task05.lean}

\newpage
\item Доказать, что для любого $\mathcal{M}\subset H$ множество $\mathcal{M}^{\perp}$ является замкнутым подпространством.

\begin{center}\small\textbf{Формализация на Lean 4}\end{center}
\inputminted{lean}{FunAn/Section12_HilbertSpaces/Task06.lean}

\newpage
\item Доказать, что если $\mathcal{M}, \mathcal{N}$ — замкнутые подпространства в $H$ и $\mathcal{M}\perp\mathcal{N}$, то $\mathcal{M}+\mathcal{N}$ также замкнутое подпространство.

\begin{center}\small\textbf{Формализация на Lean 4}\end{center}
\inputminted{lean}{FunAn/Section12_HilbertSpaces/Task07.lean}

\newpage
\item Доказать, что
$$\mathcal{M}=\left\{x\in l^{2}:x=\left(\xi_{1},\frac{\xi_{1}}{1},\xi_{3},\frac{\xi_{3}}{3},\xi_{5},\frac{\xi_{5}}{5},\dots\right)\right\},$$
и
$$\mathcal{N}=\left\{x\in l^{2}:x=(\xi_{1},0,\xi_{3},0,\xi_{5},0,\dots)\right\}$$
являются замкнутыми подпространствами, но $\mathcal{M}+\mathcal{N}$ не является замкнутым подпространством в $l^{2}$.\\
\textit{Указание.} Проверить, что $\overline{\mathcal{M}+\mathcal{N}}=l^{2}$, $\mathcal{M}+\mathcal{N}\ne l^{2}$.

\begin{center}\small\textbf{Формализация на Lean 4}\end{center}
\inputminted{lean}{FunAn/Section12_HilbertSpaces/Task08.lean}

\newpage
\item Пусть $\{e_{k}\}$ — ортонормированная система в $H$. Положим $x_{k}=(\cos\frac{1}{k})e_{2k}+(\sin\frac{1}{k})e_{2k+1}$. Доказать, что множество $L+\mathcal{M}$ не замкнуто, где $L$ — замкнутое подпространство, порожденное системой $\{e_{2k}\}$, a $\mathcal{M}$ — замкнутое подпространство, порожденное системой $\{x_{k}\}$.

\begin{center}\small\textbf{Формализация на Lean 4}\end{center}
\inputminted{lean}{FunAn/Section12_HilbertSpaces/Task09.lean}

\item Для функции $e^{t}$ найти многочлены $p_{n}(t)$ степени $n=0,1,2$ такие, что $||e^{t}-p_{n}(t)||$ минимальна в $L^{2}(-1,1)$.
\item Для функции $t^{3}$ найти многочлены $p_{n}(t)$ степени $n=0,1,2$ такие, что $||t^{3}-p_{n}(t)||$ минимальна в $L^{2}(-1,1)$.
\newpage
\item Доказать, что результатом процесса ортогонализации системы элементов $\{t^{n}\}_{n=0}^{\infty}$ в пространстве $L^{2}(-1,1)$ является система многочленов Лежандра $p_{n}(t)=c_{n}\frac{d^{n}}{dt^{n}}[(t^{2}-1)^{n}]$. Чему равны коэффициенты $c_{n}$?

\begin{center}\small\textbf{Формализация на Lean 4}\end{center}
\inputminted{lean}{FunAn/Section12_HilbertSpaces/Task12.lean}

\newpage
\item Пусть $x_{n}(t)=(-1)^{k}$, $t\in\left(\frac{k}{2^{n}},\frac{k+1}{2^{n}}\right)$, $n=0,1,\dots; k=0,1,\dots,2^{n}-1$, в концах этих интервалов $x_{n}(t)=0$. Показать, что эта система (система Радемахера) ортонормальна в $L^{2}(0,1)$, но не является базисом.\\
\textit{Указание.} Проверить, что функция $x(t)=t(t-1)+\frac{1}{6}$ ортогональна всем $x_{n}(t)$ в $L^{2}(0,1)$.

\begin{center}\small\textbf{Формализация на Lean 4}\end{center}
\inputminted{lean}{FunAn/Section12_HilbertSpaces/Task13.lean}

\newpage
\item Проверить, что система $\left\{\sqrt{\frac{2}{\pi}}\sin nt\right\}_{n=1}^{\infty}$ — ортонормальный базис в $L^{2}(0,\pi)$, но она не нормирована в $L^{2}(-\pi,\pi)$ и не является базисом.

\begin{center}\small\textbf{Формализация на Lean 4}\end{center}
\inputminted{lean}{FunAn/Section12_HilbertSpaces/Task14.lean}

\newpage
\item Доказать, что система $\{x_{n}(t)=\sin \mu_{n}t\}$, где $\mu_{n}$ — положительные корни уравнения $\operatorname{tg} \mu=\mu$, ортогональна, не нормирована и не образует базис в $L^{2}(0,1)$.

\begin{center}\small\textbf{Формализация на Lean 4}\end{center}
\inputminted{lean}{FunAn/Section12_HilbertSpaces/Task15.lean}

\item Найти $\mathcal{M}^{\perp}$ в $H$, если:
\begin{enumerate}[label=\alph*)]
\item $H=l^{2}$, $\mathcal{M}=\{(1,1,0,0,\dots)\}$;
\item $H=L^{2}(-\pi,\pi)$, $\mathcal{M}=\{\sin nt\}_{n=1}^{\infty}$;
\item $H=L^{2}(-\pi,\pi)$, $\mathcal{M}=\{\cos nt\}_{n=0}^{\infty}$.
\end{enumerate}
\newpage
\item Пусть $\mathcal{M}=\left\{x=(\xi_{1},\xi_{2},\dots)\in l^{2} : \sum_{i=1}^{\infty}\xi_{i}=0\right\}$. Доказать, что $\mathcal{M}$ — линейное многообразие, всюду плотное в $l^{2}$, то есть $\mathcal{M}^{\perp}=\{\theta\}$.

\begin{center}\small\textbf{Формализация на Lean 4}\end{center}
\inputminted{lean}{FunAn/Section12_HilbertSpaces/Task17.lean}

\end{enumerate}

\section{ЛИНЕЙНЫЕ НЕПРЕРЫВНЫЕ ФУНКЦИОНАЛЫ}

\subsection*{КОНТРОЛЬНЫЕ ВОПРОСЫ}
\begin{enumerate}
\item Дать определения однородного, аддитивного, линейного функционалов.
\item Дать определения ограниченного, непрерывного функционалов. Какова связь между непрерывностью и ограниченностью линейного функционала?
\item Привести различные определения нормы линейного непрерывного функционала.
\end{enumerate}

\subsection*{ЗАДАЧИ}
\begin{enumerate}
\item Являются ли следующие функционалы линейными, непрерывными, ограниченными в пространствах $C[0,1]$ и $L^{2}(0,1)$? Найти нормы функционалов, если они линейны и непрерывны.
\begin{enumerate}[label=\alph*)]
\item $F(x) = \int_{0}^{1} x(t) \cdot \sin t \, dt$;
\item $F(x) = \int_{0}^{1} x(t) \cdot \operatorname{sign}(t - \frac{1}{2}) \, dt$;
\item $F(x) = x(\frac{1}{2})$;
\item $F(x) = \int_{0}^{1} x(t^{2}) \sqrt{t} \, dt$;
\item $F(x) = \int_{0}^{1} x^{2}(t) \, dt$;
\item $F(x) = \int_{0}^{1} |x(t)| \, dt$;
\item $F(x) = \max_{0 \le t \le 1} x(t)$;
\item $F(x) = \int_{0}^{1/2} x(t) \, dt - \int_{1/2}^{1} x(t) \, dt$.
\end{enumerate}

\item Найти нормы функционалов, если они линейны и непрерывны:
\begin{enumerate}[label=\alph*)]
\item $F(x) = \xi_{1}$, $F: l^{2} \rightarrow \mathbb{C}$;
\item $F(x) = \sum_{k=1}^{\infty} \frac{\xi_{k}}{2^{k}}$, $F: l^{3} \rightarrow \mathbb{C}$;
\item $F(x) = \sum_{k=1}^{\infty} \frac{\xi_{k}}{k}$, $F: l^{1} \rightarrow \mathbb{C}$;
\item $F(x) = \sum_{k=1}^{\infty} \frac{\xi_{k}}{k}$, $F: l^{2} \rightarrow \mathbb{C}$;
\item $F(x) = \sum_{k=1}^{\infty} |\xi_{k}|^{3}$, $F: l^{3} \rightarrow \mathbb{C}$;
\item $F(x) = \sum_{k=1}^{\infty} (-1)^{k} \cdot \frac{\xi_{k}}{k}$, $F: l^{3} \rightarrow \mathbb{C}$.
\end{enumerate}

\item Вычислить норму функционала
$$F(x) = \int_{a}^{b} p(t)x(t) \, dt, \quad F: C[a,b] \rightarrow \mathbb{R}^{1},$$
где $p(t) \in C[a,b]$ — фиксированная функция.\\
\textit{Указание.} Воспользоваться тем, что непрерывная функция есть равномерный предел кусочно-постоянных функций.

\newpage
\item Доказать, что $F(x) = x^{\prime}(0) + x(0)$ является линейным непрерывным функционалом в пространстве $C^{(1)}[0,2]$, если $||x|| = \max_{0 \le t \le 2} |x(t)| + \max_{0 \le t \le 2} |x^{\prime}(t)|$, но не является непрерывным, если $||x|| = \max_{0 \le t \le 2} |x(t)|$.

\begin{center}\small\textbf{Формализация на Lean 4}\end{center}
\inputminted{lean}{FunAn/Section13_LinearFunctionals/Task04.lean}

\end{enumerate}

\section{ЛИНЕЙНЫЕ НЕПРЕРЫВНЫЕ ОПЕРАТОРЫ}

\subsection*{КОНТРОЛЬНЫЕ ВОПРОСЫ}
\begin{enumerate}
\item Дать определение линейного оператора.
\item Дать определения непрерывного и ограниченного оператора. Как связаны понятия непрерывности и ограниченности для линейного оператора?
\item Привести различные определения нормы линейного непрерывного оператора.
\end{enumerate}

\subsection*{ЗАДАЧИ}
\begin{enumerate}
\newpage
\item Найти общий вид линейного оператора $A: \mathbb{R}^{n} \rightarrow \mathbb{R}^{n}$. Показать, что он ограничен.


\begin{center}\small\textbf{Формализация на Lean 4}\end{center}
\inputminted{lean}{FunAn/Section14_LinearOperators/Task01.lean}

\item Проверить линейность, ограниченность, непрерывность следующих операторов. Найти их нормы, если они линейны и ограничены.
\begin{enumerate}[label=\alph*)]
\item $(Ax)(t) = \int_{0}^{t} x(\xi) \, d\xi$, $A: C[0,1] \rightarrow C[0,1]$;
\item $(Ax)(t) = \int_{0}^{1} e^{t-s} x(s) \, ds$, $A: C[0,1] \rightarrow C[2,4]$;
\item $A_{\lambda}: L^{2}(0,1) \rightarrow L^{2}(0,1)$; $(A_{\lambda}x)(t) =
\begin{cases} x(t), & \text{если } t \le \lambda; \\ 0, & \text{если } t > \lambda;
\end{cases}$
\item $(Ax)(t) = x^{\prime}(t)$, $A: C^{(1)}[-\pi,\pi] \rightarrow C[-\pi,\pi]$;
\item $(Ax)(t) = x^{\prime}(t)$, $A: C^{(1)}[-\pi,\pi] \rightarrow C[-\pi,\pi]$, $||x||_{C^{(1)}} = ||x||_{C}$;
\item $Ax = (\xi_{3}, \xi_{4}, \dots)$, $A: l^{3} \rightarrow l^{3}$;
\item $Ax = (\xi_{1}+3, \xi_{2}, \xi_{3}, \dots)$, $A: l^{\infty} \rightarrow l^{\infty}$;
\item $(Ax)(t) = \int_{8}^{t} s^{2} x^{2}(s) \, ds$, $A: C[8,10] \rightarrow C[8,10]$.
\end{enumerate}

\item Пусть $A: l^{2} \rightarrow l^{2}$, $Ax = (\alpha_{1}\xi_{1}, \alpha_{2}\xi_{2}, \dots)$, где $\alpha = \{\alpha_{k}\}_{k=1}^{\infty}$ — фиксированная числовая последовательность. Установить критерий непрерывности оператора $A$. Найти норму $A$.

\item Для каких $a(t)$ оператор $(Ax)(t) = a(t)x(t)$ непрерывен в $C[0,1]$? В $L^{2}(0,1)$? Найти нормы.

\item Для каких $\alpha, \beta$ оператор $(Ax)(t) = t^{\beta} \cdot x(t^{\alpha})$ ограничен в $L^{2}(0,1)$? Найти $||A||$.

\newpage
\item Доказать, что если $A: X \rightarrow Y$ — линейный непрерывный оператор, то он остается непрерывным при замене нормы в $X$ и $Y$ на эквивалентные.


\begin{center}\small\textbf{Формализация на Lean 4}\end{center}
\inputminted{lean}{FunAn/Section14_LinearOperators/Task06.lean}

\newpage
\item Пусть $L$ — замкнутое подпространство гильбертова пространства $H$. Доказать линейность, непрерывность и найти норму оператора $P: H \rightarrow L$, который каждому элементу $x \in H$ сопоставляет его проекцию на $L$ ($P$ называется проектором на $L$).

\begin{center}\small\textbf{Формализация на Lean 4}\end{center}
\inputminted{lean}{FunAn/Section14_LinearOperators/Task07.lean}

\end{enumerate}

\section{ОБРАТНЫЕ ОПЕРАТОРЫ}

\subsection*{КОНТРОЛЬНЫЕ ВОПРОСЫ}
\begin{enumerate}
\item Дать определение обратного оператора.
\item Сформулировать теоремы (достаточные условия) существования операторов $A^{-1}$, $(I+A)^{-1}$, $(A+\Delta A)^{-1}$.
\item Сформулировать теорему Банаха об обратном операторе.
\end{enumerate}

\subsection*{ЗАДАЧИ}
\begin{enumerate}
\newpage
\item Пусть оператор $A: \mathbb{R}^{3} \rightarrow \mathbb{R}^{3}$ определяется равенством:
$$
\begin{pmatrix} \eta_{1} \\ \eta_{2} \\ \eta_{3}
\end{pmatrix} =
\begin{pmatrix} 5 & 2 & 3 \\ -1 & 3 & -2 \\ 5 & 4 & 2
\end{pmatrix} \cdot
\begin{pmatrix} \xi_{1} \\ \xi_{2} \\ \xi_{3}
\end{pmatrix} $$
Доказать ограниченность оператора $A$, оценить $||A||$, найти $A^{-1}$.


\begin{center}\small\textbf{Формализация на Lean 4}\end{center}
\inputminted{lean}{FunAn/Section15_InverseOperators/Task01.lean}

\item При каких $\alpha_{1}, \alpha_{2}, \dots$ существует обратный оператор к оператору $Ax = (\alpha_{1}\xi_{1}, \alpha_{2}\xi_{2}, \dots)$, $A: l^{2} \rightarrow l^{2}$? Когда $A^{-1}$ непрерывен?

\newpage
\item Проверить, что оператор дифференцирования $\frac{d}{dt}: C^{(1)}[0,1] \rightarrow C[0,1]$ имеет правый обратный, но не имеет левого обратного.


\begin{center}\small\textbf{Формализация на Lean 4}\end{center}
\inputminted{lean}{FunAn/Section15_InverseOperators/Task03.lean}

\item Найти обратный оператор к оператору $\frac{d}{dt}: X \rightarrow C[0,1]$, где $X = \{x \in C^{(1)}[0,1] : x(0) = 0\}$.

\item Пусть $X = \{x \in C^{(2)}[0,1] : x(0) = x(1) = 0\}$. Найти обратный к оператору $\frac{d^{2}}{dt^{2}} + \lambda: X \rightarrow C[0,1]$, если он существует. Для каких $\lambda$ существует обратный оператор?

\item При каких $\lambda$ существует обратный оператор к оператору
$$(Ax)(t) = \int_{0}^{t} x(\xi) \, d\xi + \lambda \cdot x(t), \quad A: C[0,1] \rightarrow C[0,1] ?$$
Построить $A^{-1}$.

\item Проверить, существует ли непрерывный обратный оператор к оператору $A: l^{2} \rightarrow l^{2}$, если
\begin{enumerate}[label=\alph*)]
\item $Ax = (0, \xi_{1}, \xi_{2}, \dots)$;
\item $Ax = (\xi_{1} + \xi_{2}, \xi_{2}, \xi_{3}, \dots)$;
\item $Ax = (\xi_{2} - \xi_{1}, \xi_{2} + \xi_{3}, 2\xi_{2} - 2\xi_{1}, \xi_{4}, \xi_{5}, \dots)$;
\item $Ax = (\xi_{3}, \xi_{1}, \xi_{2}, \xi_{4}, \xi_{5}, \dots)$.
\end{enumerate}

\item Пусть $(Ax)(t) = (t+1)x(t)$, $(Bx)(t) = x(t^{2})$ — операторы в $C[0,1]$. Чему равны $(AB)^{-1}, (BA)^{-1}$?

\item Пусть $A, B \in \mathcal{L}(X,X)$. Доказать, что
\begin{enumerate}[label=\alph*)]
\item если существуют $A^{-1}, B^{-1}$, то $(AB)^{-1} = B^{-1}A^{-1}$;
\newpage
\item если существуют $A^{-1}, (BA)^{-1}$, то существует $B^{-1}$.
\end{enumerate}
\textit{Указание.} $B^{-1} = A(BA)^{-1}$.


\begin{center}\small\textbf{Формализация на Lean 4}\end{center}
\inputminted{lean}{FunAn/Section15_InverseOperators/Task09.lean}

\newpage
\item Пусть $X$ — банахово пространство с нормами $|| \cdot ||_{1}$ и $|| \cdot ||_{2}$. Если $\exists C_{1} > 0$ такое, что $C_{1}||x||_{1} \le ||x||_{2}$ для $x \in X$, то $\exists C_{2} < \infty$ такое, что $||x||_{2} \le C_{2}||x||_{1}$, $x \in X$. Доказать.


\begin{center}\small\textbf{Формализация на Lean 4}\end{center}
\inputminted{lean}{FunAn/Section15_InverseOperators/Task10.lean}

\newpage
\item Пусть $A, B \in \mathcal{L}(X,X)$, причем существует $(I - AB)^{-1}$. Доказать, что существует $(I - BA)^{-1}$.\\
\textit{Указание.} $(I - BA)^{-1} = I + B(I - AB)^{-1}A$.

\begin{center}\small\textbf{Формализация на Lean 4}\end{center}
\inputminted{lean}{FunAn/Section15_InverseOperators/Task11.lean}

\end{enumerate}

\section{СОПРЯЖЕННЫЕ ПРОСТРАНСТВА. РЕФЛЕКСИВНОСТЬ}

\subsection*{КОНТРОЛЬНЫЕ ВОПРОСЫ}
\begin{enumerate}
\item Дать определение сопряженного пространства.
\item Каков общий вид линейных непрерывных функционалов в пространствах $\mathbb{R}^{n}, H, C[a,b], l^{p}, L^{p}(a,b)$, где $1 < p < \infty$?
\item Описать сопряженные пространства к пространствам $\mathbb{R}^{n}, H, l^{p}, L^{p}(a,b)$, где $1 < p < \infty$.
\item Какое пространство называется вторым сопряженным? Описать естественное отображение $\pi: X \rightarrow X^{**}$ и дать определение рефлексивного пространства $X$. Привести примеры рефлексивных пространств.
\end{enumerate}

\subsection*{ЗАДАЧИ}
\begin{enumerate}
\item Описать сопряженное пространство к пространству $\mathbb{R}^{n}$ с нормой
\begin{enumerate}[label=\alph*)]
\item $||x|| = \max_{1 \le k \le n} |\xi_{k}|$;
\item $||x|| = \sum_{k=1}^{n} |\xi_{k}|$;
\item $||x|| = (\sum_{k=1}^{n} |\xi_{k}|^{p})^{1/p}$, где $1 < p < \infty$, $x = (\xi_{1}, \dots, \xi_{n})$.
\end{enumerate}
Является ли $\mathbb{R}^{n}$ рефлексивным?

\item Описать сопряженное пространство к пространству $l^{1}$.

\item Описать сопряженное пространство к пространству $c_{0}$ сходящихся к нулю числовых последовательностей с нормой $||x|| = \max_{k \ge 1} |\xi_{k}|$. Рефлексивно ли пространство $c_{0}$?

\item Описать сопряженное пространство к пространству $c$ сходящихся к конечному пределу последовательностей с нормой $||x|| = \sup_{k \ge 1} |\xi_{k}|$. Является ли пространство $c$ рефлексивным?

\item Найти общий вид линейных непрерывных функционалов в пространстве $s$.
\end{enumerate}

\section{ТЕОРЕМА ХАНА-БАНАХА. СИЛЬНАЯ И СЛАБАЯ СХОДИМОСТИ}

\subsection*{КОНТРОЛЬНЫЕ ВОПРОСЫ}
\begin{enumerate}
\item Сформулировать теорему Хана-Банаха и следствия из нее.
\item Дать определения сильной и слабой сходимости последовательности функционалов, определение слабой сходимости последовательности элементов.
\item Сформулировать критерий слабой сходимости последовательности функционалов и элементов.
\end{enumerate}

\subsection*{ЗАДАЧИ}
\begin{enumerate}
\item Найти все линейные непрерывные продолжения функционала $f(x) = 4\xi_{1}$, $x = \xi_{1} \in \mathbb{R}^{1}$ на пространство $\mathbb{R}^{3}$ с сохранением нормы, если $\mathbb{R}^{3}$ рассматривается с нормой
\begin{enumerate}[label=\alph*)]
\item $||x|| = (\xi_{1}^{2} + \xi_{2}^{2} + \xi_{3}^{2})^{1/2}$;
\item $||x|| = |\xi_{1}| + |\xi_{2}| + |\xi_{3}|$;
\item $||x|| = \max(|\xi_{1}|, |\xi_{2}|, |\xi_{3}|)$.
\end{enumerate}

\item Найти все линейные непрерывные продолжения функционала $f(x) = \xi_{1} + \xi_{2}$, $x \in \mathbb{R}^{2}$, на пространство $l^{p}$ с сохранением нормы, если $||x|| = (|\xi_{1}|^{p} + |\xi_{2}|^{p})^{1/p}$:
\begin{enumerate}[label=\alph*)]
\item $p=1$;
\item $1 < p < \infty$.
\end{enumerate}

\newpage
\item Доказать, что последовательность функционалов
$$F_{n}(x) = \int_{-\pi}^{\pi} x(t) e^{int} \, dt$$
слабо сходится в $L^{2}(-\pi,\pi)$, но сильной сходимости нет.


\begin{center}\small\textbf{Формализация на Lean 4}\end{center}
\inputminted{lean}{FunAn/Section17_HahnBanach/Task03.lean}

\newpage
\item Покажите, что существует последовательность конечных линейных комбинаций функционалов $F_{t}(x) = x(t)$, сходящаяся слабо к функционалу $F(x) = \int_{0}^{1} x(t) \, dt$ (функционалы определены в пространстве $C[0,1]$). Имеется ли последовательность, сходящаяся по норме?


\begin{center}\small\textbf{Формализация на Lean 4}\end{center}
\inputminted{lean}{FunAn/Section17_HahnBanach/Task04.lean}

\item Пусть $F_{n}(x) = n \cdot \left[x(\frac{1}{n}) + x(-\frac{1}{n}) - 2x(0)\right]$. Показать, что последовательность $\{F_{n}\}$:
\begin{enumerate}[label=\alph*)]
\item сходится по норме к $F(x) \equiv 0$ в $C^{(2)}[-1,1]$;
\item сходится слабо к $0$ в $C^{(1)}[-1,1]$, но не сходится сильно;
\newpage
\item не сходится даже слабо в $C[-1,1]$.
\end{enumerate}
\textit{Указания.} (i) Использовать формулу Тейлора. (ii) Оценить нормы $||F_{n}||$ снизу:
$$x_{n}^{\prime}(t) =
\begin{cases} 1, & \text{если } t > \frac{1}{n}, \\ nt, & \text{если } |t| \le \frac{1}{n}, \\ -1, & \text{если } t < -\frac{1}{n}.
\end{cases}$$
(iii) Показать, что последовательность $\{||F_{n}||\}$ неограничена.


\begin{center}\small\textbf{Формализация на Lean 4}\end{center}
\inputminted{lean}{FunAn/Section17_HahnBanach/Task05.lean}

\item Пусть $F_{n}(x) = n^{2} \cdot \left[x(\frac{1}{n}) + x(-\frac{1}{n}) - 2x(0)\right]$, $F(x) = x^{\prime\prime}(0)$. Показать, что последовательность $\{F_{n}\}$:
\begin{enumerate}[label=\alph*)]
\item сходится к $F$ по норме в $C^{(3)}[-1,1]$;
\item сходится слабо к $F$ в $C^{(2)}[-1,1]$, но не сходится сильно;
\newpage
\item не сходится даже слабо в $C^{(1)}[-1,1]$.
\end{enumerate}


\begin{center}\small\textbf{Формализация на Lean 4}\end{center}
\inputminted{lean}{FunAn/Section17_HahnBanach/Task06.lean}

\newpage
\item Доказать, что слабая сходимость в $l^{2}$ сильнее покоординатной сходимости.


\begin{center}\small\textbf{Формализация на Lean 4}\end{center}
\inputminted{lean}{FunAn/Section17_HahnBanach/Task07.lean}

\item Какие из последовательностей сходятся сильно, слабо в $l^{2}$?
\begin{enumerate}[label=\alph*)]
\item $x_{n} = (1, \frac{1}{2}, \dots, \frac{1}{n}, 0, \dots)$;
\item $x_{n} = (\underbrace{0, \dots, 0}_{n-1}, 1, 0, \dots)$;
\item $x_{n} = (0, 0, \dots, 0, 1, \frac{1}{2}, \frac{1}{3}, \dots)$ (начало с $n$-го места);
\item $x_{n} = (0, \dots, 0, \frac{1}{n+1}, \frac{1}{n+2}, \dots)$;
\item $x_{n} = (\underbrace{1, \dots, 1}_{n}, \frac{1}{n}, \frac{1}{n+1}, \dots)$;
\item $x_{n} = (\frac{1}{n}, \frac{1}{n+1}, \dots)$.
\end{enumerate}

\item Исследовать на сильную и слабую сходимости следующие последовательности в $L^{2}(0,1)$:
\begin{enumerate}[label=\alph*)]
\item $x_{n}(t) = t^{n}$;
\item $x_{n}(t) =
\begin{cases} \sqrt{n}, & \text{если } t \in [0, \frac{1}{n}], \\ 0, & \text{если } t \in (\frac{1}{n}, 1];
\end{cases}$
\item $x_{n}(t) = e^{i\sqrt{n}t}$;
\item $x_{n}(t) = \sin(t \ln n)$;
\item $x_{n}(t) = \cos(n^{2}t)$;
\item $x_{n}(t) = \sin(e^{n}t)$;
\item $x_{n}(t) =
\begin{cases} \sqrt{n} \cdot (1 - nt), & \text{если } t \in [0, \frac{1}{n}), \\ 0, & \text{если } t \in [\frac{1}{n}, 1];
\end{cases}$
\item $x_{n}(t) = t^{n} - t^{n+1}$;
\item $x_{n}(t) =
\begin{cases} 2n(1 - nt), & \text{если } t \in [0, \frac{1}{n}), \\ 0, & \text{если } t \in [\frac{1}{n}, 1].
\end{cases}$
\end{enumerate}

\item Пусть последовательность $\{x_{n}\} \subset H$ слабо сходится к $x$. Доказать, что $\{x_{n}\}$ будет сходиться сильно при выполнении одного из условий:
\begin{enumerate}[label=\alph*)]
\item $||x_{n}|| \rightarrow ||x||$;
\item $||x_{n}|| \le ||x||$ $\forall n > N_{0}$;
\newpage
\item $\lim_{n \rightarrow \infty} ||x_{n}|| \le ||x||$.
\end{enumerate}

\begin{center}\small\textbf{Формализация на Lean 4}\end{center}
\inputminted{lean}{FunAn/Section17_HahnBanach/Task10.lean}

\end{enumerate}

\section{СОПРЯЖЕННЫЕ ОПЕРАТОРЫ}

\subsection*{КОНТРОЛЬНЫЕ ВОПРОСЫ}
\begin{enumerate}
\item Дать определения сопряженного оператора к линейному непрерывному оператору в нормированном и гильбертовом пространствах.
\item Какая связь между сопряженным оператором в $H$ и сопряженным оператором в нормированном пространстве, когда в качестве нормированного пространства рассматривается $H$?
\item Перечислить свойства сопряженного оператора.
\end{enumerate}

\subsection*{ЗАДАЧИ}
\begin{enumerate}
\item Найти сопряженный оператор к оператору $A: \mathbb{R}^{n} \rightarrow \mathbb{R}^{m}$, определяемому равенством
$$
\begin{pmatrix} \eta_{1} \\ \eta_{2} \\ \dots \\ \eta_{m}
\end{pmatrix} =
\begin{pmatrix} a_{11} & a_{12} & \dots & a_{1n} \\ a_{21} & a_{22} & \dots & a_{2n} \\ \dots & \dots & \dots & \dots \\ a_{m1} & a_{m2} & \dots & a_{mn}
\end{pmatrix}
\begin{pmatrix} \xi_{1} \\ \xi_{2} \\ \dots \\ \xi_{n}
\end{pmatrix} $$

\item Найти сопряженный оператор к проектору $P: H \rightarrow L$, где $L$ — замкнутое подпространство в $H$.

\item Найти сопряженный оператор к оператору $A: X \rightarrow X$, если
\begin{enumerate}[label=\alph*)]
\item $X = l^{2}$, $Ax = (0, \xi_{1}, \xi_{2}, \dots)$;
\item $X = l^{2}$, $Ax = (\alpha_{1}\xi_{1}, \alpha_{2}\xi_{2}, \dots)$;
\item $X = l^{1}$, $Ax = (\xi_{1}, \dots, \xi_{n}, 0, 0, \dots)$;
\item $X = l^{p}$, $Ax = (\underbrace{0, \dots, 0}_{n}, \xi_{1}, 0, 0, \dots)$, $1 \le p < \infty$;
\item $X = l^{p}$, $Ax = (\xi_{2}, \xi_{1}, \xi_{4}, \xi_{3}, \dots)$;
\item $X = L^{2}(0,1)$, $(Ax)(t) = \int_{0}^{t} x(\xi) \, d\xi$;
\item $X = L^{2}(0,1)$, $(Ax)(t) = x(t^{\alpha})$;
\item $X = L^{2}(0,1)$, $(A_{\lambda}x)(t) =
\begin{cases} x(t), & \text{если } 0 \le t \le \lambda, \\ 0, & \text{если } \lambda < t \le 1
\end{cases}$, \\ где $\lambda \in (0,1)$ — фиксированное число;
\item $X = H$, $Ax = (x, y) \cdot z$, где $y, z \in H$ фиксированы;
\item $X = L^{2}(0,1)$, $(Ax)(t) = \int_{0}^{1} t \cdot x(s) \, ds$.
\end{enumerate}

\newpage
\item Пусть оператор $\Phi: \mathcal{L}(X,Y) \rightarrow \mathcal{L}(Y^{*},X^{*})$ определяется равенством $\Phi(A) = A^{*}$, где $X, Y$ — нормированные пространства. Доказать, что $\Phi$ непрерывен.


\begin{center}\small\textbf{Формализация на Lean 4}\end{center}
\inputminted{lean}{FunAn/Section18_AdjointOperators/Task04.lean}

\newpage
\item Доказать, что если $A \in \mathcal{L}(X,Y)$ и $A^{-1} \in \mathcal{L}(Y,X)$, то существует $(A^{*})^{-1} \in \mathcal{L}(X^{*},Y^{*})$, причём $(A^{*})^{-1} = (A^{-1})^{*}$.


\begin{center}\small\textbf{Формализация на Lean 4}\end{center}
\inputminted{lean}{FunAn/Section18_AdjointOperators/Task05.lean}

\newpage
\item Пусть $X$ — рефлексивное банахово пространство, $Y$ — нормированное, $A \in \mathcal{L}(X,Y)$. Доказать, что $(A^{*})^{*} = A$ как оператор из $X$ в $Y$.

\begin{center}\small\textbf{Формализация на Lean 4}\end{center}
\inputminted{lean}{FunAn/Section18_AdjointOperators/Task06.lean}

\end{enumerate}

\section{РАВНОМЕРНАЯ, СИЛЬНАЯ, СЛАБАЯ СХОДИМОСТИ ПОСЛЕДОВАТЕЛЬНОСТЕЙ ОПЕРАТОРОВ}

\subsection*{КОНТРОЛЬНЫЕ ВОПРОСЫ}
\begin{enumerate}
\item Дать определения равномерной, сильной и слабой сходимостей последовательности линейных непрерывных операторов.
\item Сформулировать принцип равномерной ограниченности Банаха-Штейнхауса.
\item Сформулировать критерий сильной сходимости линейных непрерывных операторов.
\end{enumerate}

\subsection*{ЗАДАЧИ}
\begin{enumerate}
\newpage
\item Доказать, что если $A_{n} \in \mathcal{L}(\mathbb{R}^{m},\mathbb{R}^{m})$, то из слабой сходимости последовательности $\{A_{n}\}$ следует равномерная сходимость.


\begin{center}\small\textbf{Формализация на Lean 4}\end{center}
\inputminted{lean}{FunAn/Section19_OperatorConvergence/Task01.lean}

\item Пусть $A_{n}x = (0, \dots, 0, \frac{\xi_{n+1}}{n+1}, \frac{\xi_{n+2}}{n+2}, \dots)$, $B_{n}x = (\xi_{1}, \dots, \xi_{n}, 0, 0, \dots)$, $C_{n}x = (\underbrace{0, 0, \dots, 0}_{n}, \xi_{1}, 0, 0, \dots)$ — операторы в $l^{2}$. Доказать, что:
\begin{enumerate}[label=\alph*)]
\item $\{A_{n}\}$ сходится равномерно к нулевому оператору;
\item $\{B_{n}\}$ сильно сходится к единичному оператору, а равномерно не сходится;
\newpage
\item $\{C_{n}\}$ слабо сходится к нулевому оператору, но сильно не сходится.
\end{enumerate}


\begin{center}\small\textbf{Формализация на Lean 4}\end{center}
\inputminted{lean}{FunAn/Section19_OperatorConvergence/Task02.lean}

\item Пусть $A_{n}x = (\xi_{n+1}, \xi_{n+2}, \dots)$, $A_{n}: l^{2} \rightarrow l^{2}$. Доказать, что
\begin{enumerate}[label=\alph*)]
\item $\{A_{n}\}$ сходится сильно, $\{A_{n}^{*}\}$ сходится слабо к нулевому оператору;
\item $\{A_{n}A_{n}^{*}\}$ не сходится слабо к нулевому оператору;
\newpage
\item $\{A_{n}^{*}A_{n}\}$ сходится к нулевому оператору сильно.
\end{enumerate}


\begin{center}\small\textbf{Формализация на Lean 4}\end{center}
\inputminted{lean}{FunAn/Section19_OperatorConvergence/Task03.lean}

\newpage
\item Пусть $A_{n}x = (\xi_{n+1}, \xi_{n+2}, \dots)$, $A_{n}: l^{2} \rightarrow l^{2}$. Доказать, что $A_{n} \rightarrow 0$ сильно. Верно ли утверждение $A_{n}^{*} \rightarrow 0$ сильно?


\begin{center}\small\textbf{Формализация на Lean 4}\end{center}
\inputminted{lean}{FunAn/Section19_OperatorConvergence/Task04.lean}

\item Исследовать последовательности $\{A_{n}\} \subset \mathcal{L}(X,X)$ на равномерную, сильную, слабую сходимость.
\begin{enumerate}[label=\alph*)]
\item $X = l^{2}$, $A_{n}x = (\underbrace{0, \dots, 0}_{n-1}, \xi_{n}, 0, 0, \dots)$;
\item $X = L^{2}(0,1)$, $(A_{n}x)(t) = \int_{0}^{1} t^{n}\xi^{n}x(\xi) \, d\xi$;
\item $X = C[0,1]$, $(A_{n}x)(t) = t^{n}x(t)$;
\item $X = C[0,1]$, $(A_{n}x)(t) = n \cdot \int_{t}^{t+\frac{1}{n}} x(\xi) \, d\xi$;
\item $X = C[0,1]$, $(A_{n}x)(t) = t^{n} \cdot (1-t) \cdot x(t)$;
\item $X = C[0,1]$, $(A_{n}x)(t) = x(t^{1+\frac{1}{n}})$.
\end{enumerate}

\newpage
\item Пусть $A, A_{n} \in \mathcal{L}(H,H)$, $\{x_{n}\} \subset H$. Доказать, что если $A_{n} \rightarrow A$ слабо, а $x_{n} \rightarrow x$ сильно, то $A_{n}x_{n} \rightarrow Ax$ слабо. Привести пример, показывающий, что из сильной сходимости $\{A_{n}\}$ к $A$ и слабой сходимости $\{x_{n}\}$ к $x$ не следует слабая сходимость $\{A_{n}x_{n}\}$ к $Ax$.\\
\textit{Указание.} См. задачу 3.

\begin{center}\small\textbf{Формализация на Lean 4}\end{center}
\inputminted{lean}{FunAn/Section19_OperatorConvergence/Task06.lean}

\end{enumerate}

\section{КОМПАКТНЫЕ МНОЖЕСТВА}

\subsection*{КОНТРОЛЬНЫЕ ВОПРОСЫ}
\begin{enumerate}
\item Дать определение компактного в себе множества, компакта.
\item Сформулировать теоремы Арцела, Рисса (критерии компактности множества в $C[a,b], L^{p}(a,b), l^{p}$).
\end{enumerate}

\subsection*{ЗАДАЧИ}
\begin{enumerate}
\item Пусть $\mathcal{M} = \{x(t)\}$ — ограниченное множество функций из $C[a,b]$. Доказать, что множество функций вида:
\begin{enumerate}[label=\alph*)]
\item $y(s) = \int_{a}^{b} (s^{2} + t^{2}s)x(t) \, dt$;
\item $y(s) = \int_{a}^{b} (3st + s^{2})x(t) \, dt$;
\newpage
\item $y(s) = \int_{a}^{b} s \cdot x(t) \, dt$
\end{enumerate}
компактны в $L^{p}(a,b)$.


\begin{center}\small\textbf{Формализация на Lean 4}\end{center}
\inputminted{lean}{FunAn/Section20_CompactSets/Task01.lean}

\item Пусть $\mathcal{M} = \{x(t)\}$ — ограниченное множество функций из $L^{p}(a,b)$. Компактны ли в пространстве $L^{p}(a,b)$ множества функций вида:
\begin{enumerate}[label=\alph*)]
\item $y(s) = \int_{a}^{b} (3s + s^{2}t^{2})x(t) \, dt$;
\item $y(t) = t^{2}x(t)$;
\item $y(t) = x(t) + 3$;
\item $y(s) = \int_{a}^{b} K(s,t)x(t) \, dt$, где функция $K(s,t)$ непрерывна в квадрате $[a,b] \times [a,b]$?
\end{enumerate}

\item Пусть $\mathcal{M} = \{x(t)\}$ — ограниченное множество пространства $C[a,b]$. Доказать, что множества функций вида:
\begin{enumerate}[label=\alph*)]
\item $y(t) = \int_{a}^{t} x(\xi) \, d\xi$;
\newpage
\item $y(s) = \int_{a}^{b} K(s,t)x(t) \, dt$, где $K(s,t) \in C([a,b] \times [a,b])$
\end{enumerate}
компактны в пространстве $C[a,b]$.


\begin{center}\small\textbf{Формализация на Lean 4}\end{center}
\inputminted{lean}{FunAn/Section20_CompactSets/Task03.lean}

\item Будет ли ограниченное множество многочленов степени не выше $n$ компактным в себе множеством в $C[a,b]$?

\newpage
\item Доказать, что всякое множество, компактное в пространстве $C^{(1)}[a,b]$, является компактным и в пространстве $C[a,b]$.


\begin{center}\small\textbf{Формализация на Lean 4}\end{center}
\inputminted{lean}{FunAn/Section20_CompactSets/Task05.lean}

\item Компактны ли следующие множества функций в пространстве $C[0,1]$:
\begin{enumerate}[label=\alph*)]
\item $x_{n}(t) = t^{n}$, $n=1,2,\dots$;
\item $x_{n}(t) = \sin nt$, $n=1,2,\dots$;
\item $x_{n}(t) = \sin(n+t)$, $n=1,2,\dots$;
\item $x_{\alpha}(t) = \sin \alpha t$, $\alpha \in \mathbb{R}^{1}$;
\item $x_{\alpha}(t) = \sin \alpha t$, $\alpha \in [1,2]$;
\item $x_{\alpha}(t) = \operatorname{arctg} \alpha t$, $\alpha \in \mathbb{R}^{1}$;
\item $x_{\alpha}(t) = e^{t-\alpha}$, $\alpha \ge 0$?
\end{enumerate}

\item Доказать, что всякое компактное множество в пространстве
\begin{enumerate}[label=\alph*)]
\item $C[a,b]$;
\newpage
\item $l^{2}$
\end{enumerate}
нигде не плотно.


\begin{center}\small\textbf{Формализация на Lean 4}\end{center}
\inputminted{lean}{FunAn/Section20_CompactSets/Task07.lean}

\newpage
\item Доказать, что параллелепипед
$$ \mathcal{P} = \{x = (\xi_{1}, \xi_{2}, \dots) \in l^{2} : |\xi_{k}| \le \frac{1}{k}, \forall k \ge 1\} $$
является компактным в себе множеством в пространстве $l^{2}$.


\begin{center}\small\textbf{Формализация на Lean 4}\end{center}
\inputminted{lean}{FunAn/Section20_CompactSets/Task08.lean}

\item При каком условии на последовательность $\lambda_{n} \in \mathbb{R}^{1}$, $\lambda_{n} > 0$, являются компактным в себе множеством в пространстве $l^{2}$:
\begin{enumerate}[label=\alph*)]
\item параллелепипед $\mathcal{P} = \{x = (\xi_{1}, \xi_{2}, \dots) \in l^{2} : |\xi_{k}| \le \lambda_{k}, \forall k \ge 1\}$;
\item эллипсоид $\mathcal{E} = \{x = (\xi_{1}, \xi_{2}, \dots) \in l^{2} : \sum_{n=1}^{\infty} \frac{|\xi_{n}|^{2}}{\lambda_{n}^{2}} \le 1\}$?
\end{enumerate}

\item Какие из следующих множеств компактны в пространстве $C[0,1]$:
\begin{enumerate}[label=\alph*)]
\item $\{y(t) = \int_{0}^{1} e^{ts}\varphi(s) \, ds : |\varphi(s)| \le 1, \varphi \in C[0,1]\}$;
\item $\{x(t) \in C^{(2)}[0,1] : |x(t)| \le B_{0}; |x^{\prime}(t)| \le B_{1}; |x^{\prime\prime}(t)| \le B_{2}\}$;
\item $\{x(t) \in C^{(1)}[0,1] : |x(t)| \le B_{0}; |x^{\prime}(t)| \le B_{1}\}$;
\item $\{x(t) \in C^{(2)}[0,1] : |x(t)| \le B_{0}; |x^{\prime\prime}(t)| \le B_{2}\}$;
\item $\{x(t) \in C^{(2)}[0,1] : |x^{\prime}(t)| \le B_{1}; |x^{\prime\prime}(t)| \le B_{2}\}$;
\item $\{x(t) \in C[0,1] : |x(t)| \le B, |x(t_{1}) - x(t_{2})| \le L \cdot |t_{1} - t_{2}|\}$;
\item $\{x(t) \in C^{(1)}[0,1] : |x(t)| \le B, |x^{\prime}(t_{1}) - x^{\prime}(t_{2})| \le L \cdot |t_{1} - t_{2}|\}$.
\end{enumerate}

\item Компактны ли следующие множества в пространстве $L^{2}(0,1)$:
\begin{enumerate}[label=\alph*)]
\item $\{t^{\alpha}\}$, если $\alpha > -\frac{1}{2}$, если $\alpha > 0$;
\item $\{\sin nt\}$, $n \ge 1$;
\item $\{x(t) = \int_{0}^{t} \varphi(\xi) \, d\xi : \int_{0}^{1} |\varphi(\xi)|^{2} \, d\xi \le 1, \varphi \in L^{2}(0,1)\}$;
\item $\{(\ln t)^{n}\}$, $n \ge 1$.
\end{enumerate}

\item Пусть $\{c_{n}\}$ — числовая последовательность. Пусть, далее
$$ \mathcal{M} = \{x = (\xi_{1}, \xi_{2}, \dots) \in l^{2} : \sum_{k=1}^{\infty} |c_{k}\xi_{k}|^{2} \le 1\} $$
Для каких последовательностей множество $\mathcal{M}$ компактно?
\end{enumerate}

\section{КОМПАКТНЫЕ ОПЕРАТОРЫ}

\subsection*{КОНТРОЛЬНЫЕ ВОПРОСЫ}
\begin{enumerate}
\item Дать определение компактного оператора.
\item Перечислить свойства компактного оператора.
\item Сформулируйте критерии компактности множеств в пространствах $\mathbb{R}^{n}, C[a,b], L^{p}(a,b), l^{p}$.
\end{enumerate}

\subsection*{ЗАДАЧИ}
\begin{enumerate}
\newpage
\item Доказать, что любой линейный оператор $A: \mathbb{R}^{n} \rightarrow \mathbb{R}^{m}$ является компактным.


\begin{center}\small\textbf{Формализация на Lean 4}\end{center}
\inputminted{lean}{FunAn/Section21_CompactOperators/Task01.lean}

\newpage
\item Проверить компактность оператора $A: l^{2} \rightarrow l^{2}$, если
$$ Ax = A(\xi_{1}, \xi_{2}, \dots) = (0, \frac{\xi_{1}}{1}, \frac{\xi_{2}}{2}, \dots, \frac{\xi_{n}}{n}, \dots) $$


\begin{center}\small\textbf{Формализация на Lean 4}\end{center}
\inputminted{lean}{FunAn/Section21_CompactOperators/Task02.lean}

\newpage
\item Пусть $H$ — гильбертово пространство и $A: H \rightarrow H$ — линейный оператор. Доказать, что $A$ компактен тогда и только тогда, когда он любую слабо сходящуюся последовательность переводит в сильно сходящуюся.


\begin{center}\small\textbf{Формализация на Lean 4}\end{center}
\inputminted{lean}{FunAn/Section21_CompactOperators/Task03.lean}

\newpage
\item Пусть $A: l^{2} \rightarrow l^{2}$, $A(\xi_{1}, \xi_{2}, \dots) = (a_{1}\xi_{1}, a_{2}\xi_{2}, \dots)$. Доказать, что оператор $A$ компактен тогда и только тогда, когда $a_{n} \rightarrow 0$ при $n \rightarrow \infty$.


\begin{center}\small\textbf{Формализация на Lean 4}\end{center}
\inputminted{lean}{FunAn/Section21_CompactOperators/Task04.lean}

\newpage
\item Доказать, что оператор $A: C[a,b] \rightarrow C[a,b]$, $(Ax)(t) = x_{0}(t)x(t)$, где $0 \ne x_{0} \in C[a,b]$, не является компактным.

\begin{center}\small\textbf{Формализация на Lean 4}\end{center}
\inputminted{lean}{FunAn/Section21_CompactOperators/Task05.lean}

\newpage
\item Доказать компактность следующих операторов:
\begin{enumerate}[label=\alph*)]
\item $A: C[0,1] \rightarrow C[0,1]$, $(Ax)(t) = \int_{0}^{t} x_{0}(t)x(s) \, ds$, где $x_{0}(t) \in C[a,b]$;
\item $A: C[a,b] \rightarrow C[a,b]$, $(Ax)(t) = \int_{a}^{b} K(s,t)x(s) \, ds$, где $K(s,t)$ — непрерывная на $[a,b] \times [a,b]$ функция.
\end{enumerate}

\begin{center}\small\textbf{Формализация на Lean 4}\end{center}
\inputminted{lean}{FunAn/Section21_CompactOperators/Task06.lean}

\newpage
\item Доказать компактность следующих операторов:
\begin{enumerate}[label=\alph*)]
\item $A: L^{2}(0,1) \rightarrow L^{2}(0,1)$, $(Ax)(t) = \int_{0}^{t} x(s) \, ds$;
\item $A: L^{2}(0,1) \rightarrow L^{2}(0,1)$, $(Ax)(t) = x_{0}(t) \cdot \int_{0}^{t} x(s) \, ds$, где $x_{0}(t) \in C[0,1]$;
\item $A: L^{2}(a,b) \rightarrow L^{2}(a,b)$, $(Ax)(t) = \int_{a}^{b} K(s,t)x(s) \, ds$, где $K(s,t)$ — непрерывная на $[a,b] \times [a,b]$ функция;
\item $A: L^{2}(a,b) \rightarrow L^{2}(a,b)$, $(Ax)(t) = \int_{a}^{b} K(s,t)x(s) \, ds$, где $K(s,t) \in L^{2}([a,b] \times [a,b])$.
\end{enumerate}

\begin{center}\small\textbf{Формализация на Lean 4}\end{center}
\inputminted{lean}{FunAn/Section21_CompactOperators/Task07.lean}

\newpage
\item Показать, что оператор $A: l^{2} \rightarrow \mathbb{R}^{1}$, $Ax = \sum_{k=1}^{\infty} \frac{\xi_{k}}{2^{k}}$ является компактным.

\begin{center}\small\textbf{Формализация на Lean 4}\end{center}
\inputminted{lean}{FunAn/Section21_CompactOperators/Task08.lean}

\newpage
\item Показать, что образ компактного оператора является сепарабельным.


\begin{center}\small\textbf{Формализация на Lean 4}\end{center}
\inputminted{lean}{FunAn/Section21_CompactOperators/Task09.lean}

\newpage
\item Являются ли компактными следующие операторы в пространствах $C[0,1]$ и $L^{2}(0,1)$?
\begin{enumerate}[label=\alph*)]
\item $(Ax)(t) = \int_{0}^{1} x(s) \, ds$;
\item $(Ax)(t) = \int_{0}^{1} x(s)(t-s)^{-2/5} \, ds$;
\item $(Ax)(t) = \int_{0}^{1} (ts + t^{2}s^{2})x(s) \, ds$.
\end{enumerate}

\newpage
\item Доказать, что оператор дифференцирования $\frac{d}{dt}$ является компактным из $C^{(2)}[0,1]$ в $C[0,1]$.

\begin{center}\small\textbf{Формализация на Lean 4}\end{center}
\inputminted{lean}{FunAn/Section21_CompactOperators/Task11.lean}

\newpage
\item При каких $\alpha, \beta, \gamma > 0$ является компактным в пространстве $C[0,1]$ оператор $(Ax)(t) = \int_{0}^{1} t^{\alpha} s^{\beta} x(s^{\gamma}) \, ds$?

\item При каких $\alpha, \beta, \gamma$ является компактным в пространстве $L^{2}(0,1)$ оператор $(Ax)(t) = \int_{0}^{1} t^{\alpha} s^{\beta} x(s^{\gamma}) \, ds$?
\end{enumerate}

\end{document}